\chapter{从零开始的理论预备}
\label{ch:primer}

本章的定位是“真题解答前的最小理论闭包”。它不是把四门核心课完整重写一遍，而是按求真书院 AI 方向考纲，把后面真题会反复调用的定义、公式、证明套路和考核方式整理成一条可学习路径。读这一章时，不要只背公式；每个概念都要同时回答三件事：

\begin{enumerate}
  \item 数学对象是什么，例如假设类、风险、值函数、注意力矩阵；
  \item 它在算法里如何表示，例如数组、参数向量、损失函数、优化器状态；
  \item 考试会如何考，例如证明界、推导梯度、解释局限、比较模型。
\end{enumerate}

\source{结构依据：求真书院 AI 方向博士生资格考试大纲 \parencite{qzc-ai-syllabus-2026}。机器学习理论主要参考 \textcite{shalev2014uml} 与 \textcite{mohri2018foundations}；深度学习参考 \textcite{goodfellow2016deep} 与 \textcite{bishop2023deep}；优化参考 \textcite{nesterov2018convex,beck2017first}；NLP 参考 \textcite{jurafsky2026slp,goldberg2017nnnlp}。}

\paragraph{陈述层级}

本章用 \definitionlabel、\theoremlabel、\proofsketchlabel、\heuristiclabel 与 \engineeringlabel 区分形式结论、证明路线、直觉与工程经验。没有标签的连续正文仅承担叙述和衔接，不应被误读为无条件定理；每个带条件的结论以其所在段落的假设为准。

\paragraph{本章符号约定}

向量默认用小写字母，例如 \(x,w,z\)；矩阵默认用大写字母，例如 \(X,Q,K\)。概率分布用花体或罗马体，例如 \(\Dcal\) 和 \(p_\theta\)；训练样本统一写为 \(S=(z_i)_{i=1}^m\)，其中 \(z_i=(x_i,y_i)\)。经验风险与真实风险分别为 \(L_S\) 与 \(L_{\Dcal}\)。当 \(y\) 同时可能表示单个标签和标签向量时，单个标签写作 \(y_i\)，标签向量写作 \(\mathbf{y}\)。指示函数统一写作 \(\mathbb{I}\{\cdot\}\)。

跨模块处采用下表的消歧约定；局部推导沿用其上下文中已经定义的符号，但不得借用同一符号表示表中两个不同对象。

\begin{longtable}{L{0.25\textwidth}L{0.20\textwidth}L{0.45\textwidth}}
\toprule
对象 & 记号 & 说明 \\
\midrule
\endfirsthead
\toprule
对象 & 记号 & 说明 \\
\midrule
\endhead
数据分布与样本 & \(\Dcal\), \(S=(z_i)_{i=1}^m\) & \(z_i=(x_i,y_i)\)，训练样本量为 \(m\) \\
VC 维与输入维数 & \(d_{\mathrm{VC}}\), \(d_{\mathrm{in}}\) & 不再用同一 \(d\) 同时指代二者 \\
模型维数与网络深度 & \(d_{\mathrm{model}}\), \(L_{\mathrm{net}}\) & Transformer 宽度与网络层数 \\
平滑常数与损失 & \(L_{\mathrm{sm}}\), \(\ell\) 或 \(\Lcal\) & \(L_{\Dcal}\) 保留给风险，避免混淆 \\
注意力 value 矩阵 & \(V_{\mathrm{attn}}\) & 与强化学习价值函数 \(V^\pi\) 区分 \\
核函数与测量数 & \(k(x,x')\), \(M\) & \(M\) 为压缩感知测量数 \\
稀疏度、参数量、token 数 & \(s\), \(N_{\mathrm{param}}\), \(N_{\mathrm{tok}}\) & 避免与样本量、状态或矩阵列数混用 \\
\bottomrule
\end{longtable}

\paragraph{考纲覆盖索引}

下面的表不是复习内容的替代品，而是阅读地图：左列是官方考纲条目，中列说明本章对应位置，右列列出主要参考来源。每个条目的正式定义、推导和考试提醒仍以对应小节正文为准。

\begingroup
\small
\begin{longtable}{L{0.34\textwidth}L{0.34\textwidth}L{0.22\textwidth}}
\toprule
考纲条目 & 本章覆盖位置 & 主要来源 \\
\midrule
\endfirsthead
\toprule
考纲条目 & 本章覆盖位置 & 主要来源 \\
\midrule
\endhead
监督学习和 PAC 学习理论 & 监督学习、PAC、可实现 ERM、agnostic PAC、统一收敛 & \cite{shalev2014uml,mohri2018foundations} \\
VC 维、Sauer 引理、样本复杂度 & VC 维、增长函数、Sauer 引理、VC 泛化界推导 & \cite{mohri2018foundations} \\
模型选择与偏差方差 & approximation/estimation error、bias-variance 分解 & \cite{shalev2014uml} \\
线性回归和逻辑回归 & 正规方程、岭回归、逻辑损失的 MLE 与梯度 & \cite{shalev2014uml} \\
决策树、最近邻 & 信息增益/Gini、过拟合、\(k\)-NN 的 bias-variance 与距离尺度 & \cite{shalev2014uml} \\
SVM 与核方法 & margin、hinge loss、SVM 对偶、kernel trick & \cite{mohri2018foundations} \\
压缩感知 & 稀疏测量、\(\ell_1\) 松弛、RIP 唯一性 & \cite{candes2006compressive} \\
凸集、凸函数、次微分 & 凸一阶条件、次梯度、Fermat rule、sum rule、\(\ell_1\) 次微分 & \cite{nesterov2018convex,beck2017first} \\
GD、SGD、随机坐标下降 & descent lemma、强凸线性收敛、SGD telescoping、坐标更新 & \cite{beck2017first,shalev2014uml} \\
动量和自适应学习率 & momentum、Nesterov、AdaGrad、Adam 的状态变量与适用场景 & \cite{nesterov2018convex,beck2017first} \\
神经网络框架、训练、逼近和泛化 & MLP、反向传播、逼近/优化/泛化三分法、归一化、残差 & \cite{goodfellow2016deep,bishop2023deep} \\
视觉任务 & 特征提取、恢复、三维重构、光流、识别与分割的输出结构 & \cite{szeliski2022vision,bishop2023deep} \\
无监督学习、隐式正则化、迁移和元学习 & autoencoder、contrastive learning、pretraining、fine-tuning、meta-objective & \cite{goodfellow2016deep,bishop2023deep} \\
视觉生成模型和表达力分析 & VAE、GAN、flow、diffusion 的目标、likelihood、可逆性和采样代价 & \cite{goodfellow2014gan,kingma2014vae,dinh2017realnvp,ho2020ddpm,song2021score} \\
强化学习 & MDP、Bellman 方程、Q-learning、policy gradient、baseline & \cite{mohri2018foundations,sutton2018rl} \\
统计建模和 Word2vec & \(n\)-gram MLE、skip-gram、negative sampling、GloVe & \cite{jurafsky2026slp,goldberg2017nnnlp} \\
Transformer self-attention & \(Q,K,V\) 维度、\(\sqrt{d_k}\) 缩放、causal mask & \cite{vaswani2017attention,jurafsky2026slp} \\
大模型预训练对齐与 Scaling Law & next-token CE、SFT、RLHF/DPO、幂律与 compute-optimal trade-off & \cite{kaplan2020scaling,hoffmann2022training,rafailov2023dpo} \\
Embedding 知识库推理 & TransE、RAG、MLN、能力边界和失败模式 & \cite{lewis2020rag,jurafsky2026slp} \\
\bottomrule
\end{longtable}
\endgroup

\section{监督学习、PAC、VC 与泛化}

\teaching
监督学习的基本问题是：从有限样本中学一个预测规则\footnote{预测规则就是一个函数 \(h:\mathcal{X}\to\mathcal{Y}\)：看到新输入 \(x\) 时，输出 \(h(x)\)。形式上，先设一对随机变量 \((X,Y)\sim\Dcal\)，其中 \(X\) 是输入、\(Y\) 是答案；“训练集由 \(\Dcal\) 独立同分布抽样”是指每个 \((x_i,y_i)\) 都按同一个联合分布 \(P_{XY}\) 生成，且不同样本彼此独立。二分类中，\(x\) 可为邮件的词频和发件人特征，\(h(x)\in\{\text{垃圾},\text{正常}\}\)；回归中，\(x\) 可为房屋面积和位置，\(h(x)\in\R\) 预测房价；概率分类中，模型先输出 \(p_\theta(y\mid x)\)，再以 \(\argmax_y p_\theta(y\mid x)\) 作为类别预测。样本只是未知总体分布的一小批随机观测，因此训练误差小并不自动说明新样本上的误差也小。}，并希望它在未见样本上也表现好。这里的核心困难不是“训练集上能不能拟合”，而是“训练集误差能不能代表总体误差”。

设输入空间为 \(\mathcal{X}\)，标签空间为 \(\mathcal{Y}\)，未知数据分布为 \(\Dcal\)。训练集

\[
  S=((x_1,y_1),\ldots,(x_m,y_m))
\]

由 \(\Dcal\) 独立同分布抽样得到。假设类 \(\Hcal\) 是候选预测器集合，损失函数 \(\ell(h,(x,y))\) 衡量预测器 \(h\) 在样本 \((x,y)\) 上的错误。经验风险和真实风险分别为

\[
  L_S(h)=\frac1m\sum_{i=1}^{m}\ell(h,(x_i,y_i)),
  \qquad
  L_{\Dcal}(h)=\E_{(x,y)\sim\Dcal}\ell(h,(x,y)).
\]

经验风险是可计算的训练误差；真实风险是我们真正关心但无法直接计算的泛化误差。经验风险最小化（ERM）写作

\[
  h_S\in \argmin_{h\in\Hcal}L_S(h).
\]

这一定义背后的直觉很简单：既然看不到总体分布，就先选训练集上表现最好的模型。但 ERM 是否可靠取决于 \(\Hcal\) 的复杂度。如果 \(\Hcal\) 太大，它可以把随机噪声也拟合掉，训练误差会系统性偏低。

\subsection{PAC 学习}

\concepttarget{concept:pac}

PAC 是 Probably Approximately Correct 的缩写。它把“学会”写成一个概率命题：样本足够多时，学习算法以高概率输出一个误差不超过目标精度的假设。

在二分类可实现情形中，假设存在目标函数 \(f\in\Hcal\)，标签由 \(y=f(x)\) 生成。算法 \(A\) PAC 学会 \(\Hcal\)，意思是对任意精度 \(\varepsilon>0\) 和失败概率 \(\delta>0\)，当样本量

\[
  m\ge m_{\Hcal}(\varepsilon,\delta)
\]

时，有

\[
  \Pbb_{S\sim \Dcal^m}
  \left[
    L_{\Dcal}(A(S))\le \varepsilon
  \right]
  \ge 1-\delta.
\]

这里 \(\varepsilon\) 是 accuracy parameter，控制允许多错；\(\delta\) 是 confidence parameter，控制允许多大概率失败。样本复杂度 \(m_{\Hcal}(\varepsilon,\delta)\) 告诉我们要多少样本才能达到这个保证。

为什么 Bernoulli 和 union bound 会不断出现？对固定的 \(h\)，二分类 \(0\)-\(1\) 损失可写成误差指示变量

\[
  Z_i=\mathbb{I}\{h(x_i)\ne f(x_i)\}.
\]

于是

\[
  L_S(h)=\frac1m\sum_{i=1}^{m}Z_i,
  \qquad
  \E[Z_i]=L_{\Dcal}(h).
\]

固定 \(h\) 时，问题就是 Bernoulli 均值集中；若要同时控制所有 \(h\in\Hcal\)，就要把坏事件

\[
  \left|L_S(h)-L_{\Dcal}(h)\right|>\varepsilon
\]

对整个假设类做并集控制。这就是有限假设类界中 \(\log|\Hcal|\) 的来源。若 \(\Hcal\) 有限，Hoeffding 不等式和 union bound 给出典型形式

\[
  \Pbb\left[
    |L_S(h)-L_{\Dcal}(h)|>\varepsilon
  \right]
  \le
  2\exp(-2m\varepsilon^2)
\]

对每个固定 \(h\) 成立。因为

\[
  \left\{
    \exists h\in\Hcal:
    |L_S(h)-L_{\Dcal}(h)|>\varepsilon
  \right\}
  =
  \bigcup_{h\in\Hcal}
  \left\{
    |L_S(h)-L_{\Dcal}(h)|>\varepsilon
  \right\},
\]

所以由 union bound，

\[
  \Pbb\left[
    \exists h\in\Hcal:
    |L_S(h)-L_{\Dcal}(h)|>\varepsilon
  \right]
  \le
  2|\Hcal|\exp(-2m\varepsilon^2).
\]

令右边不超过 \(\delta\)，得到

\[
  m
  \ge
  \frac{1}{2\varepsilon^2}
  \log\frac{2|\Hcal|}{\delta}.
\]

这条公式不是为了背常数，而是为了看清结构：样本量随 \(1/\varepsilon^2\) 增长，随 \(\log(1/\delta)\) 增长，随模型类大小的对数增长。

\paragraph{可实现 ERM 的 PAC 推导}

考试中更常见的是可实现情形：算法输出任意训练误差为零的 \(h_S\)。要证明 \(h_S\) 真误差小，只需排除“坏但碰巧一致”的假设。称 \(h\) 是坏假设，如果

\[
  L_{\Dcal}(h)>\varepsilon.
\]

对固定坏假设 \(h\)，一个样本不暴露它错误的概率至多为 \(1-\varepsilon\)。由于样本独立，\(m\) 个样本都不暴露它错误的概率至多为

\[
  (1-\varepsilon)^m
  \le
  \exp(-m\varepsilon).
\]

若 \(\Hcal\) 有限，存在某个坏假设仍与训练集完全一致的概率不超过

\[
  |\Hcal|\exp(-m\varepsilon).
\]

令该概率不超过 \(\delta\)，得到可实现 PAC 的样本量条件

\[
  m
  \ge
  \frac{1}{\varepsilon}
  \log\frac{|\Hcal|}{\delta}.
\]

这条界和上面的二侧统一收敛界常数与 \(\varepsilon\) 阶不同，原因是它用了更强的“训练误差为零”和“可实现”假设 \parencite{shalev2014uml,mohri2018foundations}。

\examnote
考题常要求从“固定 \(h\)”过渡到“存在 \(h\)”：第一步写 Bernoulli 误差变量，第二步用集中不等式，第三步对 \(\Hcal\) 做 union bound。若漏掉第三步，通常只能得到单个假设的界。

\subsection{Agnostic PAC 与统一收敛}

\concepttarget{concept:agnostic-pac}

现实中目标函数未必属于 \(\Hcal\)，甚至标签可能含噪声。Agnostic PAC 不假设可实现性，而是要求算法输出接近类内最优的假设：

\[
  L_{\Dcal}(A(S))
  \le
  \inf_{h\in\Hcal}L_{\Dcal}(h)+\varepsilon
\]

以至少 \(1-\delta\) 的概率成立。这里 \(\inf_{h\in\Hcal}L_{\Dcal}(h)\) 是\emph{类内最优风险}（class-optimal risk），并不是 approximation error。相对于 Bayes 风险 \(L^*_{\mathrm{Bayes}}\) 的近似误差应写作 \(\inf_{h\in\Hcal}L_{\Dcal}(h)-L^*_{\mathrm{Bayes}}\)。下文若使用近似 ERM，还须把优化误差单列；本段先分析精确 ERM。ERM 的困难变成证明统一收敛：

\[
  \sup_{h\in\Hcal}
  |L_S(h)-L_{\Dcal}(h)|
  \le \varepsilon.
\]

一旦统一收敛成立，ERM 的泛化分析几乎是机械的。令 \(h^*\in\argmin_{h\in\Hcal}L_{\Dcal}(h)\)，则

\[
  L_{\Dcal}(h_S)
  \le L_S(h_S)+\varepsilon
  \le L_S(h^*)+\varepsilon
  \le L_{\Dcal}(h^*)+2\varepsilon.
\]

第一步和第三步来自统一收敛，第二步来自 ERM 定义。这类“三行证明”是学习理论真题最常见的骨架。

如果题目要求最终 excess risk 不超过 \(\varepsilon\)，通常把统一收敛精度设成 \(\varepsilon/2\)。即若

\[
  \sup_{h\in\Hcal}
  |L_S(h)-L_{\Dcal}(h)|
  \le \frac{\varepsilon}{2},
\]

则同样推导给出

\[
  L_{\Dcal}(h_S)
  \le
  L_{\Dcal}(h^*)+\varepsilon.
\]

这就是 agnostic ERM 样本复杂度里经常出现额外常数的来源，而不是新的理论困难。

\subsection{VC 维、增长函数与 Sauer 引理}

\concepttarget{concept:vc-sauer}

当 \(\Hcal\) 无限时，不能再用 \(|\Hcal|\) 直接度量复杂度。VC 理论度量的是假设类在有限点集上实现任意标记的能力。

\definitionlabel 给定 \(m\) 个点 \(x_1,\ldots,x_m\)，增长函数定义为

\[
  \tau_{\Hcal}(m)
  =
  \max_{x_1,\ldots,x_m}
  \left|
  \{
    (h(x_1),\ldots,h(x_m)):h\in\Hcal
  \}
  \right|.
\]

如果\emph{存在}一个 \(m\) 点集，其所有 \(2^m\) 种二值标记都能由 \(\Hcal\) 实现，就说该点集被 \(\Hcal\) shatter。用上确界而不是最大值，VC 维的严格定义为

\[
  d_{\mathrm{VC}}(\Hcal)
  =
  \sup\{m\in\mathbb{N}:\tau_{\Hcal}(m)=2^m\},
\]

Sauer 引理说明有限 VC 维会把指数级增长压成多项式增长：若 \(\operatorname{VCdim}(\Hcal)=d<\infty\)，则

\[
  \tau_{\Hcal}(m)
  \le
  \sum_{i=0}^{d}\binom{m}{i}
  \le
  \left(\frac{em}{d}\right)^d
  \quad (m\ge d).
\]

并约定集合无上界时 \(d_{\mathrm{VC}}(\Hcal)=\infty\)。这条引理的意义是：即使 \(\Hcal\) 无限，只要 VC 维有限，在 \(m\) 个样本上它能产生的标记模式也只有多项式多；但要得到随机样本上的统一收敛，不能直接把有限类证明中的 \(|\Hcal|\) 替换成 \(\tau_{\Hcal}(m)\)。

\paragraph{从 Sauer 到 VC 泛化界}

\proofsketchlabel 以下骨架针对二值类的 \(0\)--\(1\) 损失（更一般地，可替换为有界损失并相应调整常数），并取 \(S,S'\overset{\mathrm{iid}}{\sim}\Dcal^m\)。令 \(S'\) 是独立 ghost sample。先把 \(L_{\Dcal}\) 与 \(L_{S'}\) 比较，再经 symmetrization 把偏差写成 \(S\cup S'\) 上的两份经验风险之差。给定这 \(2m\) 个点后，\(\Hcal\) 只会产生至多 \(\tau_{\Hcal}(2m)\) 个标记模式，最后才对这些模式使用 union bound。因此在常数与满足 symmetrization 前提的样本量范围内，有

\[
  \Pbb\left[
    \exists h\in\Hcal:
    |L_S(h)-L_{\Dcal}(h)|>\varepsilon
  \right]
  \lesssim
  \tau_{\Hcal}(2m)\exp(-c m\varepsilon^2),
\]

其中 \(c>0\) 是集中不等式给出的常数。若 \(\operatorname{VCdim}(\Hcal)=d\)，Sauer 引理给出

\[
  \log \tau_{\Hcal}(2m)
  \le
  d\log\left(\frac{2em}{d}\right).
\]

令失败概率不超过 \(\delta\)，需要

\[
  d\log\left(\frac{2em}{d}\right)
  -
  c m\varepsilon^2
  \le
  \log\delta.
\]

移项后可读出

\[
  \varepsilon
  \gtrsim
  \sqrt{
    \frac{
      d\log(2em/d)+\log(1/\delta)
    }{m}
  }.
\]

这就是 VC 泛化界的来源。考试若要求“说明样本复杂度和 VC 维的关系”，应写出 ghost sample、symmetrization、\(\tau_{\Hcal}(2m)\) 与 Sauer 引理这条链，而不是把增长函数直接代入固定假设的 Hoeffding 界。

典型的 VC 泛化界写作

\[
  L_{\Dcal}(h)
  \le
  L_S(h)
  +
  O\left(
    \sqrt{\frac{d\log(m/d)+\log(1/\delta)}{m}}
  \right)
\]

对所有 \(h\in\Hcal\) 同时成立。考试中不必执着于常数，但必须看懂 \(d/m\) 是主导量：样本量远大于 VC 维时，训练误差才会可靠。

\subsection{Rademacher complexity}

Rademacher complexity 是另一种更细的复杂度度量。它问的是：函数类能多好地拟合随机噪声？给定样本 \(S=(z_1,\ldots,z_m)\)，经验 Rademacher complexity 为

\[
  \widehat{\mathfrak{R}}_S(\mathcal{G})
  =
  \E_{\sigma}
  \left[
    \sup_{g\in\mathcal{G}}
    \frac1m
    \sum_{i=1}^{m}\sigma_i g(z_i)
  \right],
\]

其中 \(\sigma_i\) 独立均匀取值于 \(\{-1,+1\}\)。如果 \(\mathcal{G}\) 能让内积 \(\sum_i\sigma_i g(z_i)\) 很大，说明它可以跟随机标签高度相关，复杂度高，泛化风险也高。它比 VC 维更数据依赖，在 margin 类和核方法中尤其自然。

\paragraph{Rademacher 界的推导入口}

设 \(\mathcal{G}\) 是取值于 \([0,1]\) 的损失函数类。我们想控制

\[
  \sup_{g\in\mathcal{G}}
  \left(
    \E g(z)-\frac1m\sum_{i=1}^{m}g(z_i)
  \right).
\]

引入一个与 \(S\) 独立的 ghost sample \(S'=(z_1',\ldots,z_m')\)。因为 \(\E g(z)=\E_{S'}m^{-1}\sum_i g(z_i')\)，先把总体期望替换为 ghost sample 平均：

\[
  \E_S
  \sup_g
  \left(
    \E g-\widehat{\E}_S g
  \right)
  \le
  \E_{S,S'}
  \sup_g
  \frac1m
  \sum_{i=1}^{m}
  \left(g(z_i')-g(z_i)\right).
\]

再用 Rademacher 符号随机交换 \(z_i\) 与 \(z_i'\)：

\[
  \E_{S,S'}
  \sup_g
  \frac1m
  \sum_i
  \left(g(z_i')-g(z_i)\right)
  =
  \E_{S,S',\sigma}
  \sup_g
  \frac1m
  \sum_i
  \sigma_i
  \left(g(z_i')-g(z_i)\right).
\]

由 supremum 的次可加性，右边不超过两个相同的 Rademacher 项之和，因此得到

\[
  \E_S
  \sup_g
  \left(
    \E g-\widehat{\E}_S g
  \right)
  \le
  2\mathfrak{R}_m(\mathcal{G}).
\]

这里 \(\mathfrak{R}_m(\mathcal{G})=\E_S[\widehat{\mathfrak{R}}_S(\mathcal{G})]\) 是经验 Rademacher complexity 对样本抽样再取期望后的总体版本，\(\widehat{\E}_S g=m^{-1}\sum_i g(z_i)\) 是样本平均。

高概率版本再配合 McDiarmid 或 Hoeffding 型浓缩不等式。这一推导叫 symmetrization，是 Rademacher complexity 真正进入泛化界的地方 \parencite{mohri2018foundations}。

\paragraph{实现映射}

在代码里，\(S\) 是数据矩阵和标签向量，\(h\) 是模型对象，\(L_S\) 是训练循环返回的平均 loss，\(L_{\Dcal}\) 只能用验证集或测试集近似。理论证明里常说“独立同分布抽样”；工程上对应数据切分、随机种子、去重、避免训练测试泄漏。

\paragraph{考核内容}

\begin{itemize}
  \item 写出 PAC 或 agnostic PAC 的精确定义，区分 \(\varepsilon\)、\(\delta\) 和样本复杂度。
  \item 从 Bernoulli 指示变量和 union bound 推有限类样本复杂度。
  \item 定义 shattering、增长函数、VC 维，并用 Sauer 引理得到统一收敛形式。
  \item 解释 approximation error、estimation error 与模型选择的关系。
  \item 读懂 Rademacher complexity 的定义，并说明它为什么度量“拟合随机噪声”的能力。
\end{itemize}

\section{模型选择、偏差方差与经典监督模型}

\source{本节对应考纲“模型选择、偏差方差、线性回归、逻辑回归、决策树、最近邻、SVM、核方法、压缩感知”。基础模型参考 \textcite{shalev2014uml}；核方法参考 \textcite{mohri2018foundations}；压缩感知参考 \textcite{candes2006compressive}。}

\subsection{模型选择与偏差方差}

\concepttarget{concept:bias-variance}

模型选择的核心不是选训练误差最低的模型，而是在表达能力和泛化可靠性之间平衡。把误差分解成两类更清楚：

\[
  L_{\Dcal}(h_S)-L_{\Dcal}(h^*_{\mathrm{Bayes}})
  =
  \underbrace{
    L_{\Dcal}(h^*_{\Hcal})-L_{\Dcal}(h^*_{\mathrm{Bayes}})
  }_{\text{approximation error}}
  +
  \underbrace{
    L_{\Dcal}(h_S)-L_{\Dcal}(h^*_{\Hcal})
  }_{\text{estimation error}},
\]

其中 \(h^*_{\Hcal}\) 是 \(\Hcal\) 内真实风险最小的预测器。模型类越大，approximation error 往往越小，但 estimation error 往往越大。模型类越小，估计更稳定，但可能欠拟合。

偏差方差分解是平方损失下的对应版本。设训练集随机，学习算法输出 \(\hat f_S(x)\)，对固定 \(x\) 有

\[
  \E_S[(\hat f_S(x)-f(x))^2]
  =
  \left(\E_S[\hat f_S(x)]-f(x)\right)^2
  +
  \E_S\left[
    \left(\hat f_S(x)-\E_S[\hat f_S(x)]\right)^2
  \right].
\]

第一项是 bias squared，表示平均预测离真函数多远；第二项是 variance，表示换一批训练集后预测波动多大。正则化、交叉验证、结构风险最小化（SRM）都可以看作在这两者之间调参。

\paragraph{偏差方差分解的推导}

令

\[
  \bar f(x)=\E_S[\hat f_S(x)].
\]

对误差项加减 \(\bar f(x)\)：

\[
  \hat f_S(x)-f(x)
  =
  \left(\hat f_S(x)-\bar f(x)\right)
  +
  \left(\bar f(x)-f(x)\right).
\]

平方并对训练集随机性取期望：

\[
  \E_S[(\hat f_S(x)-f(x))^2]
  =
  \E_S[(\hat f_S(x)-\bar f(x))^2]
  +
  (\bar f(x)-f(x))^2
  +
  2(\bar f(x)-f(x))
  \E_S[\hat f_S(x)-\bar f(x)].
\]

最后一项为零，因为

\[
  \E_S[\hat f_S(x)-\bar f(x)]
  =
  \E_S[\hat f_S(x)]-\bar f(x)
  =
  0.
\]

所以得到 bias squared 加 variance。若观测标签还含噪声 \(y=f(x)+\xi\)，且 \(\E[\xi\mid x]=0\)、\(\Var(\xi\mid x)=\sigma^2\)，预测 \(y\) 的均方误差还会多出不可约噪声 \(\sigma^2\)。

\subsection{线性回归}

线性回归假设

\[
  h_w(x)=\langle w,x\rangle+b,
\]

常用平方损失

\[
  \ell(h_w,(x,y))=(h_w(x)-y)^2.
\]

把样本写成设计矩阵 \(X\in\R^{m\times d}\)，其中第 \(i\) 行是 \(x_i^\top\)；标签向量写成 \(\mathbf{y}=(y_1,\ldots,y_m)^\top\in\R^m\)。忽略偏置或把偏置并入一列常数特征，最小二乘问题为

\[
  \min_{w\in\R^d}
  \frac1m\|Xw-\mathbf{y}\|_2^2.
\]

一阶条件给出正规方程

\[
  X^\top Xw=X^\top \mathbf{y}.
\]

若 \(X^\top X\) 可逆，则

\[
  w=(X^\top X)^{-1}X^\top \mathbf{y}.
\]

若不可逆，说明特征共线或样本不足，需要伪逆或正则化。岭回归把目标改成

\[
  \min_w
  \frac1m\|Xw-\mathbf{y}\|_2^2+\lambda\|w\|_2^2,
\]

一阶条件变成

\[
  (X^\top X+m\lambda I)w=X^\top \mathbf{y}.
\]

这解释了正则化的数值意义：它把病态矩阵往正定方向推，提高稳定性。
这里 \(I\in\R^{d\times d}\) 是单位矩阵，\(\lambda\ge0\) 是正则化强度。

\paragraph{正规方程的推导}

设

\[
  F(w)=\frac1m\|Xw-\mathbf{y}\|_2^2
  =
  \frac1m(Xw-\mathbf{y})^\top(Xw-\mathbf{y}).
\]

展开得到

\[
  F(w)
  =
  \frac1m
  \left(
    w^\top X^\top Xw
    -
    2\mathbf{y}^\top Xw
    +
    \mathbf{y}^\top\mathbf{y}
  \right).
\]

对 \(w\) 求梯度：

\[
  \nabla F(w)
  =
  \frac{2}{m}X^\top Xw
  -
  \frac{2}{m}X^\top\mathbf{y}.
\]

令 \(\nabla F(w)=0\)，即得到 \(X^\top Xw=X^\top\mathbf{y}\)。岭回归只是在梯度中多出 \(2\lambda w\)，所以

\[
  \frac{2}{m}X^\top Xw
  -
  \frac{2}{m}X^\top\mathbf{y}
  +
  2\lambda w
  =
  0,
\]

等价于

\[
  (X^\top X+m\lambda I)w=X^\top\mathbf{y}.
\]

\subsection{逻辑回归}

逻辑回归用于二分类，但输出的是概率：

\[
  p_w(y=1\mid x)=\sigma(\langle w,x\rangle+b),
  \qquad
  \sigma(t)=\frac{1}{1+\exp(-t)}.
\]

若标签 \(y\in\{-1,+1\}\)，逻辑损失为

\[
  \ell(w,(x,y))
  =
  \log\left(1+\exp(-y\langle w,x\rangle)\right).
\]

这可以从极大似然推出。它的优点是 convex 且可微，因此梯度下降能稳定求解。梯度为

\[
  \nabla_w\ell(w,(x,y))
  =
  -\frac{y x}{1+\exp(y\langle w,x\rangle)}.
\]

考试常把逻辑回归和 SVM 比较：逻辑损失平滑、概率解释强；hinge loss 更直接追求 margin；二者都是凸 surrogate loss，用可优化的损失替代 \(0\)-\(1\) 损失。

\paragraph{逻辑损失的 MLE 推导}

先用 \(y\in\{0,1\}\) 表示标签。设

\[
  p_w(y=1\mid x)=\sigma(\langle w,x\rangle+b),
  \qquad
  p_w(y=0\mid x)=1-\sigma(\langle w,x\rangle+b).
\]

单样本似然可以合并写成

\[
  p_w(y\mid x)
  =
  p^y(1-p)^{1-y},
  \qquad
  p=\sigma(\langle w,x\rangle+b).
\]

负对数似然为

\[
  -\log p_w(y\mid x)
  =
  -y\log p-(1-y)\log(1-p).
\]

若改用 \(y\in\{-1,+1\}\)，令 \(s=y(\langle w,x\rangle+b)\)。当 \(y=1\) 时，负对数似然是

\[
  -\log\sigma(\langle w,x\rangle+b)
  =
  \log(1+\exp(-s)).
\]

当 \(y=-1\) 时，正确类别概率为

\[
  1-\sigma(\langle w,x\rangle+b)
  =
  \sigma(-\langle w,x\rangle-b),
\]

负对数似然同样是 \(\log(1+\exp(-s))\)。因此统一得到

\[
  \ell(w,(x,y))
  =
  \log(1+\exp(-y(\langle w,x\rangle+b))).
\]

再令 \(s=y\langle w,x\rangle\) 并暂时忽略 \(b\)，链式法则给出

\[
  \nabla_w \ell
  =
  \frac{1}{1+\exp(-s)}
  \exp(-s)(-y x)
  =
  -\frac{y x}{1+\exp(s)}.
\]

这就是前面梯度公式。

\subsection{决策树}

决策树把输入空间递归切分成若干区域，每个叶子给一个预测。它的优势是可解释，缺点是若不限制深度或叶子数，容易过拟合。一个有 \(k\) 个叶子的树大致能记住 \(k\) 个局部区域，因此复杂度随树大小增长。

常见分裂准则是信息增益或 Gini impurity。若节点中正类比例为 \(p\)，熵为

\[
  H(p)=-p\log p-(1-p)\log(1-p),
\]

Gini impurity 为

\[
  G(p)=2p(1-p).
\]

分裂的目标是让子节点更纯。对候选分裂 \(a\)，信息增益可写成

\[
  \operatorname{Gain}(a)
  =
  H(S)
  -
  \sum_{v}
  \frac{|S_v|}{|S|}
  H(S_v).
\]

树模型考核通常不要求手推复杂证明，而是要求你说明过拟合来源、剪枝或 MDL/正则化思想，以及为什么贪心分裂不保证全局最优。

\subsection{最近邻算法}

\(k\)-nearest neighbors 不显式训练参数，而是在预测时寻找距离最近的训练样本。分类规则为

\[
  \hat y(x)
  =
  \operatorname{majority}
  \{y_i:x_i\in N_k(x)\},
\]

其中 \(N_k(x)\) 是距离 \(x\) 最近的 \(k\) 个训练点。\(k\) 小时方差大，容易被噪声影响；\(k\) 大时偏差大，局部结构被平均掉。距离度量和特征缩放极其重要，因为欧氏距离会被量纲大的特征支配。

\subsection{SVM、margin 与核方法}

\concepttarget{concept:svm-kernel}

线性 SVM 的核心是 margin。给定二分类样本 \(y_i\in\{-1,+1\}\)，硬间隔 SVM 求

\[
  \min_{w,b}\frac12\|w\|_2^2
  \quad
  \text{s.t.}
  \quad
  y_i(\langle w,x_i\rangle+b)\ge 1,\quad i=1,\ldots,m.
\]

为什么最小化 \(\|w\|\) 等价于最大化 margin？因为点到超平面的距离为

\[
  \frac{|\langle w,x\rangle+b|}{\|w\|_2}.
\]

约束把函数间隔固定为至少 \(1\)，于是几何间隔至少为 \(1/\|w\|_2\)。最小化 \(\|w\|\) 就是最大化间隔。

非可分数据用软间隔或 hinge loss：

\[
  \min_{w,b}
  \lambda\|w\|_2^2
  +
  \frac1m
  \sum_{i=1}^{m}
  \max\{0,1-y_i(\langle w,x_i\rangle+b)\}.
\]

核方法把输入映射到高维特征空间 \(\Phi(x)\)，但不显式计算 \(\Phi\)，只计算

\[
  K(x,x')=\langle \Phi(x),\Phi(x')\rangle.
\]

严格地说，实值核 \(k:\mathcal{X}\times\mathcal{X}\to\R\) 必须对称，即 \(k(x,x')=k(x',x)\)，并对任意有限点集 \(x_1,\ldots,x_m\) 以及任意系数 \(c\in\R^m\) 满足

\[
  \sum_{i,j=1}^{m}c_i c_j k(x_i,x_j)\ge 0.
\]

等价地，每一个有限 Gram 矩阵

\[
  [K(x_i,x_j)]_{i,j=1}^{m}
\]

半正定。在这一条件下，存在某个 Hilbert 空间及特征映射 \(\Phi\)，使 \(k(x,x')=\langle\Phi(x),\Phi(x')\rangle\)。这不是“任意相似度函数都可用”的许可；Mercer 型连续性与紧性假设只在需要积分算子谱展开时才额外出现。核技巧的考试重点是两点：第一，非线性边界可以变成高维空间线性分隔；第二，泛化不单由高维维数决定，margin 与正则化才是关键。

\paragraph{SVM 对偶与 kernel trick 推导}

硬间隔 SVM 的拉格朗日函数为

\[
  \mathcal{L}(w,b,\alpha)
  =
  \frac12\|w\|_2^2
  -
  \sum_{i=1}^{m}
  \alpha_i
  \left(
    y_i(\langle w,x_i\rangle+b)-1
  \right),
\]

其中 \(\alpha_i\ge0\) 是第 \(i\) 个约束的拉格朗日乘子。对 \(w\) 和 \(b\) 求驻点条件：

\[
  \nabla_w\mathcal{L}
  =
  w-\sum_{i=1}^{m}\alpha_i y_i x_i
  =
  0,
  \qquad
  \frac{\partial\mathcal{L}}{\partial b}
  =
  -\sum_{i=1}^{m}\alpha_i y_i
  =
  0.
\]

因此

\[
  w=\sum_{i=1}^{m}\alpha_i y_i x_i,
  \qquad
  \sum_{i=1}^{m}\alpha_i y_i=0.
\]

把 \(w\) 代回拉格朗日函数，得到对偶问题

\[
  \max_{\alpha\in\R^m}
  \sum_{i=1}^{m}\alpha_i
  -
  \frac12
  \sum_{i=1}^{m}
  \sum_{j=1}^{m}
  \alpha_i\alpha_j y_i y_j
  \langle x_i,x_j\rangle
\]

约束为

\[
  \alpha_i\ge0,
  \qquad
  \sum_{i=1}^{m}\alpha_i y_i=0.
\]

注意训练目标只通过内积 \(\langle x_i,x_j\rangle\) 依赖样本。若把样本替换为特征映射 \(\Phi(x_i)\)，内积变成

\[
  \langle \Phi(x_i),\Phi(x_j)\rangle=K(x_i,x_j).
\]

于是无需显式计算高维 \(\Phi(x)\)，只要在对偶目标中把 \(\langle x_i,x_j\rangle\) 替换为 \(K(x_i,x_j)\)。这就是 kernel trick。预测时

\[
  h(x)
  =
  \operatorname{sign}
  \left(
    \sum_{i=1}^{m}
    \alpha_i y_i K(x_i,x)
    +
    b
  \right).
\]

只有 \(\alpha_i>0\) 的训练点会出现在预测式中，它们就是 support vectors。

\subsection{压缩感知}

\concepttarget{concept:compressed-sensing}

压缩感知回答一个看似不可能的问题：如果未知信号 \(x\in\R^N\) 维度很高，但只有 \(s\) 个非零或近似非零分量，能否用远少于 \(N\) 个线性测量恢复它？

测量模型为

\[
  y=Mx,
  \qquad
  M\in\R^{K\times N},
  \qquad
  K\ll N.
\]

直接求最稀疏解是

\[
  \min_z \|z\|_0
  \quad
  \text{s.t.}
  \quad
  Mz=y,
\]

但这是组合优化。压缩感知的核心结果是，在随机测量或满足 null-space property，或满足具体定理所要求的足够强 RIP 条件的测量矩阵下，可以用凸松弛

\[
  \min_z \|z\|_1
  \quad
  \text{s.t.}
  \quad
  Mz=y
\]

精确或稳定恢复稀疏信号。RIP 的形式为：对所有 \(s\)-稀疏向量 \(z\)，

\[
  (1-\delta_S)\|z\|_2^2
  \le
  \|Mz\|_2^2
  \le
  (1+\delta_S)\|z\|_2^2.
\]

这表示 \(M\) 在所有稀疏子空间上近似保持距离，因此不会把两个不同的稀疏向量压到同一个测量。考纲中把压缩感知放在机器学习理论内，通常是考“稀疏性 + 凸松弛 + 随机测量 + 样本/测量复杂度”的基本思想，而不是完整调和分析。

\paragraph{RIP 为什么给唯一性}

设 \(z_1,z_2\) 都是 \(s\)-稀疏向量，且有同样测量：

\[
  Mz_1=Mz_2.
\]

令 \(u=z_1-z_2\)。两个 \(s\)-稀疏向量之差至多是 \(2s\)-稀疏，因此若 \(M\) 满足 \(2s\)-RIP，且 \(\delta_{2s}<1\)，则

\[
  (1-\delta_{2s})\|u\|_2^2
  \le
  \|Mu\|_2^2.
\]

但 \(Mu=Mz_1-Mz_2=0\)，于是

\[
  (1-\delta_{2s})\|u\|_2^2\le 0.
\]

由于 \(1-\delta_{2s}>0\)，只能有 \(u=0\)，即 \(z_1=z_2\)。所以 \(2s\)-RIP 至少保证了 \(\ell_0\) 问题中 \(s\)-稀疏解的可识别性；它本身并不自动推出 \(\ell_1\) 解的恢复。后者通常由 null-space property，或足够强的 RIP 条件推出。带噪观测 \(y=Mx+e\)、\(\|e\|_2\le\eta\) 时，BPDN

\[
  \min_z\|z\|_1\quad\text{s.t.}\quad\|Mz-y\|_2\le\eta
\]

在相应 RIP 条件下满足 \(\|\hat x-x\|_2\le C_0\eta+C_1\sigma_s(x)_1/\sqrt{s}\)，其中 \(\sigma_s(x)_1=\inf_{\|z\|_0\le s}\|x-z\|_1\)。这才是“稳定且近似稀疏”的精确定义 \parencite{candes2006compressive}。

\paragraph{考核内容}

\begin{itemize}
  \item 用 approximation error 和 estimation error 解释模型选择。
  \item 推线性回归正规方程，并说明不可逆时为什么需要正则化。
  \item 推逻辑回归损失和梯度，比较 logistic loss 与 hinge loss。
  \item 解释决策树为何过拟合，以及剪枝、深度限制、MDL 的作用。
  \item 写出 SVM hard/soft margin 目标，解释 margin 与 \(\|w\|\) 的关系。
  \item 定义核函数和 Gram 矩阵半正定条件。
  \item 对压缩感知说明 \(K\ll N\)、稀疏性、\(\ell_1\) 松弛和 RIP 的关系。
\end{itemize}

\section{优化方法：从凸性到现代一阶算法}

\source{本节对应考纲“凸集和凸函数、次微分、梯度下降、随机梯度下降、随机坐标下降、动量加速、自适应学习率”。基础参考 \textcite{nesterov2018convex,beck2017first,shalev2014uml}。}

\subsection{凸集、凸函数与一阶条件}

\concepttarget{concept:convex-first-order}

集合 \(C\subseteq\R^d\) 是凸的，表示任意 \(x,y\in C\) 和 \(\theta\in[0,1]\)，都有

\[
  \theta x+(1-\theta)y\in C.
\]

函数 \(f:C\to\R\) 是凸的，表示

\[
  f(\theta x+(1-\theta)y)
  \le
  \theta f(x)+(1-\theta)f(y).
\]

如果 \(f\) 可微，凸性等价于一阶下界

\[
  f(y)\ge f(x)+\langle \nabla f(x),y-x\rangle.
\]

这条式子的几何意义是：凸函数图像永远在任一点切平面之上。因此若 \(\nabla f(x^*)=0\)，就有

\[
  f(y)\ge f(x^*)
\]

对所有 \(y\) 成立，即任意 stationary point 都是全局最优点。这是凸优化和非凸深度学习优化的关键区别。

\paragraph{一阶条件的推导}

设 \(f\) 可微且凸。由凸性，对任意 \(t\in(0,1]\)，有

\[
  f((1-t)x+ty)
  \le
  (1-t)f(x)+tf(y).
\]

整理得

\[
  \frac{f(x+t(y-x))-f(x)}{t}
  \le
  f(y)-f(x).
\]

令 \(t\downarrow0\)，左边收敛到方向导数

\[
  \langle\nabla f(x),y-x\rangle.
\]

因此

\[
  f(y)\ge f(x)+\langle\nabla f(x),y-x\rangle.
\]

这说明凸函数的任意切平面都是全局下界。考试中若要求证明“局部最优就是全局最优”，通常就是把这个一阶下界用于 \(x=x^*\)。

若函数不可微，使用次梯度。向量 \(g\in\R^d\) 是 \(f\) 在 \(x\) 的次梯度，若

\[
  f(y)\ge f(x)+\langle g,y-x\rangle
  \quad
  \forall y.
\]

所有次梯度构成次微分 \(\partial f(x)\)。最优性条件变成

\[
  0\in \partial f(x^*).
\]

这叫 Fermat rule，是不可微凸优化里最常考的最优性条件。还需要记住两个常用性质：若 \(f\) 和 \(g\) 是适当的凸函数，并且常见正则条件满足，则

\[
  \partial(f+g)(x)
  =
  \partial f(x)+\partial g(x).
\]

若 \(A\in\R^{m\times d}\)，\(b\in\R^m\)，则仿射复合满足

\[
  \partial (f(Ax+b))
  =
  A^\top \partial f(Ax+b).
\]

这些性质解释了为什么 Lasso、hinge loss 和复合优化能用“光滑梯度 + 非光滑次梯度/prox”的形式处理 \parencite{nesterov2018convex,beck2017first}。

例如 \(f(x)=|x|\) 在 \(0\) 处不可微，但

\[
  \partial |0|=[-1,1],
\]

所以 \(0\in\partial |0|\)，\(0\) 是最小点。

\subsection{Smoothness、强凸性与下降引理}

\concepttarget{concept:smooth-strong-convex}

函数 \(f\) 是 \(L\)-smooth，表示梯度 Lipschitz：

\[
  \|\nabla f(x)-\nabla f(y)\|
  \le
  L\|x-y\|.
\]

等价地，

\[
  f(y)
  \le
  f(x)+\langle\nabla f(x),y-x\rangle
  +\frac{L}{2}\|y-x\|^2.
\]

这叫 descent lemma。把 \(y=x-\eta\nabla f(x)\) 代入，得到

\[
  f(x-\eta\nabla f(x))
  \le
  f(x)
  -
  \eta\left(1-\frac{L\eta}{2}\right)
  \|\nabla f(x)\|^2.
\]

只要 \(0<\eta<2/L\)，函数值下降。考试里常用它证明梯度下降收敛。

这一步的完整代入如下：

\[
  f(x-\eta\nabla f(x))
  \le
  f(x)
  +
  \langle\nabla f(x),-\eta\nabla f(x)\rangle
  +
  \frac{L}{2}\|-\eta\nabla f(x)\|^2.
\]

内积项为 \(-\eta\|\nabla f(x)\|^2\)，二次项为 \((L\eta^2/2)\|\nabla f(x)\|^2\)，合并后得到

\[
  f(x-\eta\nabla f(x))
  \le
  f(x)
  -
  \eta\left(1-\frac{L\eta}{2}\right)\|\nabla f(x)\|^2.
\]

函数 \(f\) 是 \(\mu\)-strongly convex，表示

\[
  f(y)\ge
  f(x)+\langle\nabla f(x),y-x\rangle
  +\frac{\mu}{2}\|y-x\|^2.
\]

强凸性至多说明\emph{若}极小点存在则它唯一；存在性还可由例如闭且有界的可行域上的下半连续性，或无约束情形中 \(f\) proper、lower-semicontinuous 且 coercive 保证。若 \(f\) 同时 \(L\)-smooth、\(\mu\)-strongly convex，且极小点 \(x^*\) 存在，梯度下降取 \(\eta=1/L\) 时有

\[
  f(x_t)-f(x^*)
  \le
  \left(1-\frac{\mu}{L}\right)^t
  \left(f(x_0)-f(x^*)\right).
\]

这里的 \(L/\mu\) 是 condition number，越大越病态，收敛越慢。

\paragraph{强凸线性收敛的推导}

取 \(x_{t+1}=x_t-\frac1L\nabla f(x_t)\)。由下降引理，

\[
  f(x_{t+1})
  \le
  f(x_t)
  -
  \frac{1}{2L}\|\nabla f(x_t)\|^2.
\]

对 \(\mu\)-strongly convex 函数，有

\[
  \|\nabla f(x)\|^2
  \ge
  2\mu(f(x)-f(x^*)).
\]

把它代入上式，得到

\[
  f(x_{t+1})-f(x^*)
  \le
  \left(1-\frac{\mu}{L}\right)
  (f(x_t)-f(x^*)).
\]

递推 \(t\) 次便得到

\[
  f(x_t)-f(x^*)
  \le
  \left(1-\frac{\mu}{L}\right)^t
  (f(x_0)-f(x^*)).
\]

这里每一步都在用同一个结构：smoothness 给下降，strong convexity 把梯度范数和函数值差联系起来。

\subsection{GD、SGD 与随机坐标下降}

\concepttarget{concept:sgd}

全批量梯度下降为

\[
  x_{t+1}=x_t-\eta_t\nabla f(x_t).
\]

若目标是经验风险

\[
  f(w)=\frac1m\sum_{i=1}^{m}\ell_i(w),
\]

则全梯度需要遍历全部样本。随机梯度下降用无偏估计

\[
  g_t=\nabla \ell_{i_t}(w_t),
  \qquad
  \E[g_t\mid w_t]=\nabla f(w_t),
\]

并更新

\[
  w_{t+1}=w_t-\eta_t g_t.
\]

SGD 的优势是单步便宜，可以在线学习；代价是方差大，需要学习率调度、mini-batch、动量或自适应方法。证明中最重要的一步是展开平方距离：

\[
  \|w_{t+1}-w^*\|^2
  =
  \|w_t-w^*\|^2
  -
  2\eta_t\langle g_t,w_t-w^*\rangle
  +
  \eta_t^2\|g_t\|^2.
\]

对条件期望取平均，再用凸性把内积下界为函数值差，即可得到收敛界。

\paragraph{SGD 收敛不等式的标准推导}

假设 \(f\) 凸，\(w^*\in\argmin_w f(w)\)，并且随机梯度满足

\[
  \E[g_t\mid w_t]=\nabla f(w_t),
  \qquad
  \E[\|g_t\|^2\mid w_t]\le G^2.
\]

从平方距离展开式出发，对 \(w_t\) 条件取期望：

\[
  \E[\|w_{t+1}-w^*\|^2\mid w_t]
  \le
  \|w_t-w^*\|^2
  -
  2\eta_t\langle\nabla f(w_t),w_t-w^*\rangle
  +
  \eta_t^2G^2.
\]

由凸性一阶条件，

\[
  f(w_t)-f(w^*)
  \le
  \langle\nabla f(w_t),w_t-w^*\rangle.
\]

于是

\[
  2\eta_t(f(w_t)-f(w^*))
  \le
  \|w_t-w^*\|^2
  -
  \E[\|w_{t+1}-w^*\|^2\mid w_t]
  +
  \eta_t^2G^2.
\]

对 \(t=0,\ldots,T-1\) 求和，距离项 telescoping：

\[
  \sum_{t=0}^{T-1}
  2\eta_t\E[f(w_t)-f(w^*)]
  \le
  \|w_0-w^*\|^2
  +
  G^2\sum_{t=0}^{T-1}\eta_t^2.
\]

若取常步长 \(\eta_t=\eta\)，并输出平均点

\[
  \bar w_T=\frac1T\sum_{t=0}^{T-1}w_t,
\]

由凸性 \(f(\bar w_T)\le T^{-1}\sum_t f(w_t)\)，得到

\[
  \E[f(\bar w_T)-f(w^*)]
  \le
  \frac{\|w_0-w^*\|^2}{2\eta T}
  +
  \frac{\eta G^2}{2}.
\]

这条推导说明 SGD 的步长选择本质是在优化误差项和噪声项之间折中。

随机坐标下降每次只更新一个坐标：

\[
  x_{t+1}
  =
  x_t-\eta_t e_{j_t}\nabla_{j_t}f(x_t).
\]

这里 \(j_t\in\{1,\ldots,d\}\) 是第 \(t\) 步随机选到的坐标，\(e_{j_t}\in\R^d\) 是第 \(j_t\) 个标准基向量，\(\nabla_{j_t}f(x_t)\) 是梯度的第 \(j_t\) 个分量。

当特征维度大、单坐标梯度便宜，或目标可分解时它很有用。考试重点是说明它与 SGD 的随机性不同：SGD 随机抽样本，坐标下降随机抽坐标。

\subsection{Proximal operator 与复合优化}

\concepttarget{concept:prox}

很多机器学习目标是光滑项加非光滑正则项：

\[
  \min_x F(x)=f(x)+R(x),
\]

例如 Lasso 中 \(R(x)=\lambda\|x\|_1\)。proximal operator 定义为

\[
  \prox_{\gamma R}(u)
  =
  \argmin_x
  \left\{
    R(x)+\frac{1}{2\gamma}\|x-u\|^2
  \right\}.
\]

prox-gradient 更新为

\[
  x_{t+1}
  =
  \prox_{\eta R}
  \left(x_t-\eta\nabla f(x_t)\right).
\]

这可以理解为：先按光滑项做一步梯度下降，再用 \(R\) 的 prox 把解拉回具有结构的区域。对 \(\ell_1\) 正则，prox 是 soft-thresholding：

\[
  [\prox_{\eta\lambda\|\cdot\|_1}(u)]_j
  =
  \operatorname{sgn}(u_j)\max\{|u_j|-\eta\lambda,0\}.
\]

因此 \(\ell_1\) 正则会产生稀疏解，不只是“让参数小”。

\paragraph{\(\ell_1\) prox 的逐坐标推导}

因为 \(\ell_1\) 正则和平方项都是逐坐标可分的，只需解一维问题

\[
  \min_x
  \lambda |x|
  +
  \frac{1}{2\eta}(x-u)^2.
\]

若 \(x>0\)，目标函数可写成

\[
  \lambda x+\frac{1}{2\eta}(x-u)^2,
\]

一阶条件为

\[
  \lambda+\frac{1}{\eta}(x-u)=0,
\]

所以 \(x=u-\eta\lambda\)。该解位于 \(x>0\) 区域，当且仅当 \(u>\eta\lambda\)。

若 \(x<0\)，目标函数可写成

\[
  -\lambda x+\frac{1}{2\eta}(x-u)^2,
\]

一阶条件为

\[
  -\lambda+\frac{1}{\eta}(x-u)=0,
\]

所以 \(x=u+\eta\lambda\)。该解位于 \(x<0\) 区域，当且仅当 \(u<-\eta\lambda\)。

若 \(|u|\le\eta\lambda\)，最优点位于不可微点 \(x=0\)。这是因为次梯度条件为

\[
  0\in
  \lambda[-1,1]
  -
  \frac{u}{\eta},
\]

等价于 \(u/\eta\in\lambda[-1,1]\)。三种情况合并得到

\[
  x^*
  =
  \operatorname{sgn}(u)\max\{|u|-\eta\lambda,0\}.
\]

\paragraph{prox 最优性、单调性与非扩张}

\concepttarget{concept:prox-nonexpansive}

考试中的 prox-SGD 递推经常不是只用 prox 的定义，而是用它的最优性条件和非扩张性。先设 \(R:\R^d\to(-\infty,+\infty]\) 是 proper closed convex 函数，步长 \(\gamma>0\)，并定义

\[
  x^+
  =
  \prox_{\gamma R}(u)
  =
  \argmin_x
  \left\{
    R(x)+\frac{1}{2\gamma}\|x-u\|^2
  \right\}.
\]

由于目标是 convex，Fermat rule 给出一阶最优性条件

\[
  0
  \in
  \partial R(x^+)
  +
  \frac1\gamma(x^+-u).
\]

移项后得到常用形式

\[
  \frac{u-x^+}{\gamma}
  \in
  \partial R(x^+).
\]

这一步是后面 fixed point 的来源。如果 \(x^*\) 是复合目标 \(F(x)=f(x)+R(x)\) 的最优点，且 \(f\) 可微、\(R\) convex，则

\[
  0\in\nabla f(x^*)+\partial R(x^*).
\]

也就是说存在 \(r^*\in\partial R(x^*)\) 使得 \(r^*=-\nabla f(x^*)\)。把 \(u=x^*-\gamma\nabla f(x^*)\) 和 \(x^+=x^*\) 代入上面的 prox 最优性条件，就得到

\[
  x^*
  =
  \prox_{\gamma R}
  \left(x^*-\gamma\nabla f(x^*)\right).
\]

接下来证明非扩张性。设

\[
  x^+=\prox_{\gamma R}(u),
  \qquad
  y^+=\prox_{\gamma R}(v).
\]

由 prox 最优性，

\[
  \frac{u-x^+}{\gamma}\in\partial R(x^+),
  \qquad
  \frac{v-y^+}{\gamma}\in\partial R(y^+).
\]

凸函数次微分的单调性说明，对任意 \(r_x\in\partial R(x)\) 和 \(r_y\in\partial R(y)\)，有

\[
  \langle r_x-r_y,x-y\rangle\ge0.
\]

把 \(x=x^+\)、\(y=y^+\)、\(r_x=(u-x^+)/\gamma\)、\(r_y=(v-y^+)/\gamma\) 代入，得到

\[
  \left\langle
    (u-x^+)-(v-y^+),
    x^+-y^+
  \right\rangle
  \ge0.
\]

展开括号：

\[
  \left\langle
    u-v,
    x^+-y^+
  \right\rangle
  -
  \|x^+-y^+\|^2
  \ge0.
\]

因此

\[
  \|x^+-y^+\|^2
  \le
  \langle u-v,x^+-y^+\rangle.
\]

这比普通非扩张更强，称为 firm nonexpansiveness。再用 Cauchy--Schwarz 不等式，

\[
  \|x^+-y^+\|^2
  \le
  \|u-v\|\,\|x^+-y^+\|.
\]

若 \(x^+\ne y^+\)，两边除以 \(\|x^+-y^+\|\)；若 \(x^+=y^+\)，结论平凡成立。于是

\[
  \|\prox_{\gamma R}(u)-\prox_{\gamma R}(v)\|
  \le
  \|u-v\|.
\]

这就是 prox 非扩张。它在考试中的作用是把含 prox 的一步递推转化为普通欧氏距离递推，例如

\[
  \|x_{k+1}-x^*\|
  \le
  \left\|
    x_k-\gamma g_k
    -
    \left(x^*-\gamma\nabla f(x^*)\right)
  \right\|.
\]

这里 \(g_k\) 是第 \(k\) 步的梯度或随机梯度估计。后续再展开平方、取条件期望，就进入 SGD 方差分解和强凸收敛分析 \parencite{beck2017first,nesterov2018convex}。

\subsection{动量、Nesterov 与自适应学习率}

动量法维护速度变量：

\[
  v_{t+1}=\beta v_t+\nabla f(x_t),
  \qquad
  x_{t+1}=x_t-\eta v_{t+1}.
\]

直觉是沿稳定下降方向累积速度，抵消高曲率方向的来回震荡。Nesterov 加速在“前瞻点”计算梯度：

\[
  v_{t+1}
  =
  \beta v_t+\nabla f(x_t-\eta\beta v_t),
  \qquad
  x_{t+1}=x_t-\eta v_{t+1}.
\]

AdaGrad、RMSProp、Adam 等方法为不同坐标设置不同有效学习率。Adam 的典型形式为

\[
  m_t=\beta_1m_{t-1}+(1-\beta_1)g_t,
  \qquad
  v_t=\beta_2v_{t-1}+(1-\beta_2)g_t\odot g_t,
\]

\[
  w_{t+1}
  =
  w_t
  -
  \eta
  \frac{\widehat m_t}{\sqrt{\widehat v_t}+\epsilon}.
\]

其中 \(\widehat m_t,\widehat v_t\) 是 bias correction 后的矩估计。考试中通常不要求证明 Adam 收敛，但要求解释为什么它能处理不同坐标尺度，以及为什么 \(\epsilon\) 防止除零。

\paragraph{考核内容}

\begin{itemize}
  \item 判断集合/函数是否凸，写出一阶凸性条件。
  \item 用次梯度定义证明不可微点的最优性。
  \item 用 smoothness 推梯度下降下降量。
  \item 比较 GD、SGD、mini-batch SGD 和随机坐标下降的计算代价与方差来源。
  \item 写出 prox 的定义，并推 \(\ell_1\) prox 的 soft-thresholding 形式。
  \item 解释动量、Nesterov、AdaGrad/Adam 的状态变量和适用场景。
\end{itemize}

\section{深度学习：网络、训练、泛化与生成模型}

\source{本节对应考纲“深度学习基础、视觉任务、无监督学习、视觉生成模型、表达力分析”。基础神经网络参考 \textcite{goodfellow2016deep,bishop2023deep,jurafsky2026slp}；视觉任务参考 \textcite{szeliski2022vision}；生成模型参考 \textcite{goodfellow2014gan,kingma2014vae,dinh2017realnvp,ho2020ddpm,song2021score}。}

\subsection{神经网络框架}

一个多层感知机（MLP）可以写成函数复合：

\[
  h_0=x,
  \qquad
  a_\ell=W_\ell h_{\ell-1}+b_\ell,
  \qquad
  h_\ell=\phi(a_\ell),
  \qquad
  \ell=1,\ldots,L.
\]

最后一层输出 logits \(o\)，分类任务通常用 softmax

\[
  p_\theta(y=k\mid x)
  =
  \frac{\exp(o_k)}{\sum_{j}\exp(o_j)}
\]

和交叉熵损失

\[
  \Lcal(\theta)
  =
  -\log p_\theta(y\mid x).
\]

神经网络的参数是所有 \(W_\ell,b_\ell\)。训练就是在样本平均损失上做随机优化。和线性模型相比，深度网络的关键变化是：特征不再固定，而是由前面层共同学习。

这里 \(h_\ell\in\R^{d_\ell}\) 是第 \(\ell\) 层激活向量，\(a_\ell\in\R^{d_\ell}\) 是激活前向量，\(W_\ell\in\R^{d_\ell\times d_{\ell-1}}\)，\(b_\ell\in\R^{d_\ell}\)，非线性函数 \(\phi\) 按坐标作用。考试推导时先写清这些维度，可以避免把 \(W_\ell^\top\) 和 \(W_\ell\) 的方向写反。

\subsection{反向传播}

\concepttarget{concept:backprop}
\concepttarget{concept:jacobian-adjoint}

反向传播只是链式法则的动态规划。设单样本损失为 \(J=\Lcal(h_L,y)\)，定义

\[
  \delta_\ell=\frac{\partial J}{\partial a_\ell}.
\]

若 \(h_\ell=\phi(a_\ell)\)，则

\[
  \delta_\ell
  =
  (W_{\ell+1}^\top \delta_{\ell+1})
  \odot
  \phi'(a_\ell),
\]

并且

\[
  \frac{\partial J}{\partial W_\ell}
  =
  \delta_\ell h_{\ell-1}^\top,
  \qquad
  \frac{\partial J}{\partial b_\ell}
  =
  \delta_\ell.
\]

这里 \(\odot\) 表示 Hadamard product，也就是按坐标相乘。

考试若要求推反向传播，不要把它神秘化：先画出计算图，确定局部 Jacobian，再把上游梯度乘下来即可。

从自动微分语言看，\(\delta_\ell\) 是中间变量 \(a_\ell\) 的 adjoint，也就是损失 \(J\) 对该变量的上游敏感度。若 \(J_\ell=\partial a_{\ell+1}/\partial a_\ell\) 是从第 \(\ell\) 层到第 \(\ell+1\) 层的局部 Jacobian，则反向传播的抽象形式是

\[
  \delta_\ell
  =
  J_\ell^\top \delta_{\ell+1}.
\]

这句话把“链式法则”翻译成可评分答案：先写局部 Jacobian，再写 adjoint 递推，最后给出各参数梯度。

\paragraph{反向传播递推的推导}

由

\[
  a_{\ell+1}=W_{\ell+1}h_\ell+b_{\ell+1},
  \qquad
  h_\ell=\phi(a_\ell),
\]

链式法则给出

\[
  \frac{\partial J}{\partial h_\ell}
  =
  W_{\ell+1}^\top
  \frac{\partial J}{\partial a_{\ell+1}}
  =
  W_{\ell+1}^\top\delta_{\ell+1}.
\]

由于 \(h_\ell=\phi(a_\ell)\) 是按坐标作用，

\[
  \frac{\partial J}{\partial a_\ell}
  =
  \frac{\partial J}{\partial h_\ell}
  \odot
  \phi'(a_\ell).
\]

合并得到

\[
  \delta_\ell
  =
  (W_{\ell+1}^\top\delta_{\ell+1})
  \odot
  \phi'(a_\ell).
\]

对权重矩阵的单个元素 \(W_{\ell,jk}\)，有

\[
  a_{\ell,j}
  =
  \sum_k W_{\ell,jk}h_{\ell-1,k}+b_{\ell,j}.
\]

因此

\[
  \frac{\partial J}{\partial W_{\ell,jk}}
  =
  \frac{\partial J}{\partial a_{\ell,j}}
  \frac{\partial a_{\ell,j}}{\partial W_{\ell,jk}}
  =
  \delta_{\ell,j}h_{\ell-1,k}.
\]

写成矩阵形式就是

\[
  \frac{\partial J}{\partial W_\ell}
  =
  \delta_\ell h_{\ell-1}^\top.
\]

sigmoid 的导数为

\[
  \sigma'(z)=\sigma(z)(1-\sigma(z))\le\frac14.
\]

深层网络反传时反复乘小因子，容易出现 vanishing gradient。ReLU

\[
  \operatorname{ReLU}(z)=\max\{z,0\}
\]

在正半轴导数为 \(1\)，缓解了系统性梯度缩小，但负半轴导数为 \(0\)，可能出现 dead ReLU。

\subsection{逼近能力、训练方法与泛化分析}

考纲把“神经网络框架、训练方法、逼近和泛化分析”放在同一条，是因为这三件事经常被混淆。逼近能力问的是“是否存在参数能表示目标函数”；训练方法问的是“优化算法能否找到好参数”；泛化分析问的是“找到的参数在未见样本上是否可靠”。三者都重要，但任何一个都不能单独推出另外两个。

经典 universal approximation 结论可以概括为：若 \(K\subset\R^d\) 是 compact set，\(f:K\to\R\) 连续，那么在适当非线性激活下，对任意 \(\varepsilon>0\)，存在足够宽的神经网络 \(F_\theta\)，使得

\[
  \sup_{x\in K}|F_\theta(x)-f(x)|<\varepsilon.
\]

这里的 \(K\) 是输入被限制的紧集，\(\theta\) 是网络参数，\(\varepsilon\) 是允许的 uniform approximation error。这条定理说明网络类表达能力很强，但它不说明 SGD 一定能找到这组参数，也不说明有限样本训练出的模型一定泛化 \parencite{goodfellow2016deep,bishop2023deep}。

训练方法通常可以抽象成随机一阶迭代：

\[
  \theta_{t+1}
  =
  \theta_t-\eta_t \widehat{\nabla}\Lcal(\theta_t),
\]

其中 \(\eta_t>0\) 是学习率，\(\widehat{\nabla}\Lcal(\theta_t)\) 是由 mini-batch、动量、归一化或自适应学习率修正后的梯度估计。考试里如果问“训练方法”，不要只写 backprop；backprop 计算梯度，SGD/Adam 这类优化器决定如何用梯度更新参数。

泛化分析则需要回到第三章开头的语言：模型类越大，approximation error 可能越小，但 estimation error 和过拟合风险可能变大。深度学习的特殊性在于，过参数化网络常常能同时插值训练数据并泛化，这通常要用隐式正则化、margin、数据增强、早停、归一化和优化偏置来解释，而不能只靠参数个数解释。

\subsection{归一化、残差与泛化}

Batch normalization 对 mini-batch 内激活做标准化：

\[
  \widehat a
  =
  \frac{a-\mu_B}{\sqrt{\sigma_B^2+\epsilon}},
  \qquad
  y=\gamma \widehat a+\beta.
\]

它的作用不是单纯“归一化数据”，而是改善优化几何，使各层输入尺度更稳定。Layer normalization 则在单个样本的特征维度上归一化，常用于 Transformer。

这里 \(\mu_B\) 和 \(\sigma_B^2\) 分别是当前 mini-batch 中某个通道或特征的均值和方差，\(\epsilon>0\) 用于数值稳定，\(\gamma\) 和 \(\beta\) 是可学习的缩放和平移参数。引入 \(\gamma,\beta\) 的原因是：标准化不应永久限制网络表达任意均值和尺度的能力。

残差连接写作

\[
  h_{\ell+1}=h_\ell+F_\ell(h_\ell).
\]

它让网络可以容易学习恒等映射，也让梯度有直接通路。Transformer、ResNet 都依赖这种结构。

深度网络泛化不能只用参数数目解释，因为过参数化网络也能泛化。常见解释包括隐式正则化、margin、数据增强、早停、归一化和优化算法偏置。考试若问“为什么深度网络参数很多仍能泛化”，不要回答成一句“正则化”，而要区分显式正则化和隐式正则化。

\subsection{视觉任务与无监督学习}

视觉任务可以按输出结构分类：

\begin{itemize}
  \item 特征提取：把图像映射为向量表示，用于检索、分类或下游任务。
  \item 图像恢复：去噪、超分辨率、去模糊，本质是从退化观测恢复干净图像。
  \item 三维重构：从多视角、深度或隐式表示恢复几何结构。
  \item 光流估计：估计相邻帧中像素或特征点的运动场。
  \item 识别和分割：分类给整图标签，检测给框，语义分割给每像素类别，实例分割区分对象实例。
\end{itemize}

考试若要求比较视觉任务，最稳妥的写法是先写输出对象，再写损失或评价指标：

\begingroup
\small
\begin{longtable}{L{0.22\textwidth}L{0.35\textwidth}L{0.32\textwidth}}
\toprule
任务 & 输出对象 & 常见训练信号 \\
\midrule
\endfirsthead
\toprule
任务 & 输出对象 & 常见训练信号 \\
\midrule
\endhead
特征提取 & 表示向量 \(z=f_\theta(x)\) & 分类损失、对比损失、检索排序损失 \\
图像恢复 & 干净图像或残差图像 \(\hat x\) & 像素损失、perceptual loss、扩散去噪损失 \\
三维重构 & 深度、点云、mesh、NeRF 或隐式场 & 重投影误差、几何一致性、体渲染损失 \\
光流估计 & 像素运动场 \(u(x)\in\R^2\) & 端点误差、亮度一致性、平滑正则 \\
识别与分割 & 类别、框、mask 或每像素标签 & 交叉熵、IoU/Dice、检测框回归损失 \\
\bottomrule
\end{longtable}
\endgroup

无监督学习不依赖人工标签。自编码器（autoencoder）学习

\[
  z=f_\phi(x),
  \qquad
  \hat x=g_\theta(z),
\]

并最小化重构误差

\[
  \|x-\hat x\|^2.
\]

对比学习则构造正负样本，使同一实例的不同增强表示接近，不同实例远离。典型损失为

\[
  \ell_i
  =
  -
  \log
  \frac{
    \exp(\operatorname{sim}(z_i,z_i^+)/\tau)
  }{
    \exp(\operatorname{sim}(z_i,z_i^+)/\tau)
    +
    \sum_{j}
    \exp(\operatorname{sim}(z_i,z_j^-)/\tau)
  }.
\]

这里 \(z_i\) 是 anchor 表示，\(z_i^+\) 是与它语义相同或来自同一实例增强的 positive 表示，\(z_j^-\) 是 negative 表示，\(\operatorname{sim}\) 通常是内积或 cosine similarity；\(\tau>0\) 是温度参数，控制 softmax 分布尖锐程度。

\paragraph{隐式正则化、迁移学习与元学习}

隐式正则化指优化过程本身偏向某些解，即使目标函数里没有显式正则项。例如同样能插值训练集的多个参数中，SGD、初始化、batch size、数据增强和归一化可能偏向 margin 更大或表示更平滑的解。它解释的是“为什么过参数化模型不一定按参数数目简单过拟合”。

迁移学习先在大数据源任务上学习表示 \(f_{\theta_0}\)，再在目标任务上微调：

\[
  \theta^*
  \in
  \argmin_\theta
  \frac1m
  \sum_{i=1}^{m}
  \ell(g_\theta(x_i),y_i),
  \qquad
  \theta\ \text{从}\ \theta_0\ \text{初始化}.
\]

这里 \(g_\theta\) 是迁移后的预测器，\(\theta_0\) 是预训练参数。元学习则把“快速适配”写进训练目标：给任务分布 \(\mathcal{T}\)，希望初始化参数 \(\theta\) 经过少量内层更新 \(U_T(\theta)\) 后在任务 \(T\) 上表现好：

\[
  \min_\theta
  \E_{T\sim\mathcal{T}}
  L_T(U_T(\theta)).
\]

这三类主题不一定要求长证明，但考试常要求区分它们的目标：隐式正则化解释泛化，迁移学习复用已有表示，元学习优化适配过程。

\subsection{生成模型}

\concepttarget{concept:generative-models}

生成模型的目标是学习数据分布 \(p_{\mathrm{data}}\)，并能采样新样本。不同模型族的差别在于如何表示 \(p_\theta(x)\) 或如何生成 \(x\)。

\paragraph{表达力与可计算性的权衡}

生成模型的“表达力分析”通常不是问参数数量，而是问模型族怎样在 likelihood、采样质量、训练稳定性和可计算性之间取舍。Flow 有显式 likelihood，但可逆性限制了网络结构；VAE 有概率图模型和 ELBO，但近似后验可能带来 posterior collapse 或模糊样本；GAN 直接学习隐式生成分布，样本锐利但训练不稳定且缺少显式 likelihood；diffusion 用很多小步去噪换来稳定训练和高质量采样，但采样成本较高。这些比较来自不同模型目标的数学形式，而不是经验标签 \parencite{goodfellow2014gan,kingma2014vae,dinh2017realnvp,ho2020ddpm,song2021score}。

\paragraph{极大似然}

若模型给出显式密度，训练常用最大似然：

\[
  \max_\theta
  \E_{x\sim p_{\mathrm{data}}}
  \log p_\theta(x).
\]

等价于最小化

\[
  \KL(p_{\mathrm{data}}\|p_\theta)
  =
  \E_{p_{\mathrm{data}}}
  \log
  \frac{p_{\mathrm{data}}(x)}{p_\theta(x)},
\]

因为 \(\E_{p_{\mathrm{data}}}\log p_{\mathrm{data}}(x)\) 与 \(\theta\) 无关。

\paragraph{VAE}

\concepttarget{concept:vae-elbo}

VAE 使用潜变量 \(z\)，生成模型为 \(p_\theta(x,z)=p(z)p_\theta(x\mid z)\)。由于真实后验 \(p_\theta(z\mid x)\) 难算，引入近似后验 \(q_\phi(z\mid x)\)。ELBO 为

\[
  \log p_\theta(x)
  \ge
  \E_{q_\phi(z\mid x)}
  \log p_\theta(x\mid z)
  -
  \KL(q_\phi(z\mid x)\|p(z)).
\]

第一项鼓励重构，第二项把编码分布拉向先验。

\paragraph{ELBO 的推导}

从边缘似然开始：

\[
  \log p_\theta(x)
  =
  \log
  \int p_\theta(x,z)\,dz.
\]

乘除同一个近似后验 \(q_\phi(z\mid x)\)：

\[
  \log p_\theta(x)
  =
  \log
  \int
  q_\phi(z\mid x)
  \frac{p_\theta(x,z)}{q_\phi(z\mid x)}
  \,dz.
\]

把积分看成关于 \(q_\phi(z\mid x)\) 的期望，并用 Jensen 不等式：

\[
  \log p_\theta(x)
  \ge
  \E_{q_\phi(z\mid x)}
  \log
  \frac{p_\theta(x,z)}{q_\phi(z\mid x)}.
\]

展开 \(p_\theta(x,z)=p(z)p_\theta(x\mid z)\)，得到

\[
\begin{aligned}
  &\E_{q_\phi(z\mid x)}\log p_\theta(x\mid z)
  +
  \E_{q_\phi(z\mid x)}\log p(z)
  -
  \E_{q_\phi(z\mid x)}\log q_\phi(z\mid x)
  \\
  &\qquad =
  \E_{q_\phi(z\mid x)}\log p_\theta(x\mid z)
  -
  \KL(q_\phi(z\mid x)\|p(z)).
\end{aligned}
\]

这就是 ELBO。考试若要求解释 VAE，必须同时说明：ELBO 是 likelihood 的下界，KL 项来自近似后验与先验之间的差距，重构项来自条件生成模型 \(p_\theta(x\mid z)\)。

重参数化技巧把

\[
  z=\mu_\phi(x)+\sigma_\phi(x)\odot\epsilon,
  \qquad
  \epsilon\sim\N(0,I)
\]

写成可反传形式。

\paragraph{GAN}

GAN 有生成器 \(G_\theta(z)\) 和判别器 \(D_\psi(x)\)。经典 minimax 目标为

\[
  \min_\theta\max_\psi
  \E_{x\sim p_{\mathrm{data}}}\log D_\psi(x)
  +
  \E_{z\sim p(z)}
  \log(1-D_\psi(G_\theta(z))).
\]

直觉是判别器学习区分真样本和生成样本，生成器学习骗过判别器。优点是样本锐利；缺点是训练不稳定、mode collapse、缺少显式 likelihood。

\paragraph{Normalizing flow}

\concepttarget{concept:flow-likelihood}

Flow 用可逆映射 \(x=f_\theta(z)\) 把简单分布 \(p_0(z)\) 推到数据空间。变量替换公式为

\[
  \log p_\theta(x)
  =
  \log p_0(z)
  -
  \log\left|\det J_{f_\theta}(z)\right|,
  \qquad
  z=f_\theta^{-1}(x).
\]

多层 flow 把 log determinant 累加。考试重点是“可逆性 + Jacobian determinant + 显式 likelihood”。

\paragraph{变量替换公式的推导}

设 \(x=f_\theta(z)\)，且 \(f_\theta\) 可逆。概率质量在变量变换下保持：

\[
  p_\theta(x)\,dx
  =
  p_0(z)\,dz.
\]

多维情形中体积元满足

\[
  dx
  =
  \left|\det J_{f_\theta}(z)\right|dz.
\]

因此

\[
  p_\theta(x)
  =
  p_0(z)
  \left|
    \det J_{f_\theta}(z)
  \right|^{-1}.
\]

取对数得到

\[
  \log p_\theta(x)
  =
  \log p_0(z)
  -
  \log\left|\det J_{f_\theta}(z)\right|.
\]

如果 \(f_\theta=f_L\circ\cdots\circ f_1\)，Jacobian determinant 由链式法则相乘，log determinant 因而相加。

\paragraph{Diffusion 与 score}

\concepttarget{concept:diffusion-score}

扩散模型先定义加噪 forward process，把 \(x_0\) 逐步变成近似高斯噪声，再学习 reverse process 去噪。离散 DDPM 常写为

\[
  q(x_t\mid x_{t-1})
  =
  \N(\sqrt{1-\beta_t}x_{t-1},\beta_t I).
\]

连续时间 VP SDE 可写作

\[
  dx=-\frac12\beta(t)x\,dt+\sqrt{\beta(t)}\,dW_t.
\]

其中 \(W_t\) 是标准 Brownian motion，\(\beta(t)>0\) 是时间 \(t\) 的噪声日程。

一般 forward SDE

\[
  dx=f(x,t)\,dt+g(t)\,dW_t
\]

的 reverse-time SDE 为

\[
  dx=
  \left[
    f(x,t)-g(t)^2\nabla_x\log p_t(x)
  \right]dt
  +
  g(t)\,d\bar W_t.
\]

这里 \(\bar W_t\) 表示反向时间下的 Brownian motion。score

\[
  \nabla_x\log p_t(x)
\]

描述当前噪声水平下密度上升最快方向。训练 score model 后，就能从噪声逐步反向采样。

离散扩散里还常用闭式边缘分布。令

\[
  \alpha_t=1-\beta_t,
  \qquad
  \bar\alpha_t=\prod_{s=1}^{t}\alpha_s.
\]

反复代入 forward process 可得

\[
  q(x_t\mid x_0)
  =
  \N(\sqrt{\bar\alpha_t}x_0,(1-\bar\alpha_t)I).
\]

因此可以直接采样

\[
  x_t
  =
  \sqrt{\bar\alpha_t}x_0
  +
  \sqrt{1-\bar\alpha_t}\epsilon,
  \qquad
  \epsilon\sim\N(0,I).
\]

DDPM 常训练网络 \(\epsilon_\theta(x_t,t)\) 预测噪声，目标近似为

\[
  \E_{t,x_0,\epsilon}
  \left[
    \|\epsilon-\epsilon_\theta(x_t,t)\|^2
  \right].
\]

这解释了为什么扩散模型的训练看起来像 denoising regression：模型学习每个噪声水平下应当去掉的噪声方向。

\paragraph{考核内容}

\begin{itemize}
  \item 写出 MLP 前向传播和反向传播递推，解释 vanishing gradient。
  \item 比较 batch normalization、layer normalization 和 residual connection 的作用。
  \item 说明表达力、优化、泛化三者不是同一个问题。
  \item 按输出结构区分视觉任务。
  \item 写出 autoencoder、contrastive learning、VAE、GAN、flow、diffusion 的核心目标。
  \item 比较四类生成模型：VAE 有 ELBO，GAN 有对抗目标，flow 有显式可逆 likelihood，diffusion 学去噪或 score。
\end{itemize}

\section{强化学习与偏好对齐}

\source{本节对应考纲“贝尔曼方程、Q-learning、Policy gradient”。基础 RL 参考 \textcite{mohri2018foundations,sutton2018rl}；LLM 偏好对齐参考 \textcite{jurafsky2026slp,rafailov2023dpo}。}

\subsection{MDP、return 与值函数}

\concepttarget{concept:mdp}

强化学习研究 agent 与 environment 的交互。Markov decision process（MDP）由状态空间 \(\mathcal{S}\)、动作空间 \(\mathcal{A}\)、转移概率 \(P(s'\mid s,a)\)、奖励 \(r(s,a)\)、折扣因子 \(\gamma\in[0,1)\) 构成。

策略 \(\pi(a\mid s)\) 给出在状态 \(s\) 下选择动作 \(a\) 的概率。折扣 return 为

\[
  G_t=\sum_{k=0}^{\infty}\gamma^k r(s_{t+k},a_{t+k}).
\]

值函数和动作值函数定义为

\[
  V^\pi(s)=\E_\pi[G_t\mid s_t=s],
  \qquad
  Q^\pi(s,a)=\E_\pi[G_t\mid s_t=s,a_t=a].
\]

值函数不是一个普通监督标签，而是对未来随机轨迹的期望。它依赖策略，因为策略决定后续会访问哪些状态。

\subsection{Bellman 方程}

\concepttarget{concept:bellman}

Bellman 方程是 return 的递归展开：

\[
  V^\pi(s)
  =
  \sum_a\pi(a\mid s)
  \left[
    r(s,a)
    +
    \gamma
    \sum_{s'}P(s'\mid s,a)V^\pi(s')
  \right].
\]

最优值函数满足 Bellman optimality equation：

\[
  V^*(s)
  =
  \max_{a}
  \left[
    r(s,a)
    +
    \gamma
    \sum_{s'}P(s'\mid s,a)V^*(s')
  \right].
\]

动作值版本为

\[
  Q^*(s,a)
  =
  r(s,a)
  +
  \gamma
  \sum_{s'}P(s'\mid s,a)
  \max_{a'}Q^*(s',a').
\]

若知道 \(Q^*\)，最优策略就是

\[
  \pi^*(s)\in\argmax_a Q^*(s,a).
\]

\paragraph{Bellman 方程的推导}

从 return 定义出发：

\[
  G_t
  =
  r(s_t,a_t)
  +
  \gamma
  \sum_{k=0}^{\infty}\gamma^k r(s_{t+1+k},a_{t+1+k})
  =
  r(s_t,a_t)+\gamma G_{t+1}.
\]

在 \(s_t=s\) 条件下取期望：

\[
  V^\pi(s)
  =
  \E_\pi[r(s_t,a_t)+\gamma G_{t+1}\mid s_t=s].
\]

先对动作 \(a\sim\pi(\cdot\mid s)\) 求和，再对下一状态 \(s'\sim P(\cdot\mid s,a)\) 求和：

\[
  V^\pi(s)
  =
  \sum_a\pi(a\mid s)
  \left[
    r(s,a)
    +
    \gamma\sum_{s'}P(s'\mid s,a)V^\pi(s')
  \right].
\]

最优 Bellman 方程只是把“按当前策略平均动作”替换成“选最优动作”。因此 expectation equation 用于评估固定策略，optimality equation 用于求最优策略。

\subsection{Q-learning}

\concepttarget{concept:q-learning}

Q-learning 是 model-free 方法，不需要知道转移概率。它用采样转移 \((s_t,a_t,r_t,s_{t+1})\) 更新

\[
  Q_{t+1}(s_t,a_t)
  =
  Q_t(s_t,a_t)
  +
  \alpha_t
  \left[
    r_t
    +
    \gamma\max_{a'}Q_t(s_{t+1},a')
    -
    Q_t(s_t,a_t)
  \right].
\]

中括号内是 temporal-difference error：

\[
  \delta_t
  =
  r_t
  +
  \gamma\max_{a'}Q_t(s_{t+1},a')
  -
  Q_t(s_t,a_t).
\]

它衡量当前 \(Q\) 估计和 Bellman backup 之间的不一致。Q-learning 是 off-policy，因为目标里用的是 \(\max_{a'}\)，不一定等于实际采样动作来自的行为策略。

这个更新式来自 Bellman optimality equation 的随机近似。若 \(Q_t\) 已接近 \(Q^*\)，则理想目标应满足

\[
  Q_t(s_t,a_t)
  \approx
  r_t+\gamma\max_{a'}Q_t(s_{t+1},a').
\]

因此把右侧视为 one-step target：

\[
  Y_t=r_t+\gamma\max_{a'}Q_t(s_{t+1},a'),
\]

再对当前估计做指数移动平均：

\[
  Q_{t+1}(s_t,a_t)
  =
  (1-\alpha_t)Q_t(s_t,a_t)+\alpha_tY_t.
\]

展开即可得到前面的 Q-learning 更新式。

\subsection{Policy gradient 与 baseline}

\concepttarget{concept:policy-gradient}

当动作空间大或策略需要可微参数化时，直接优化策略更自然。设

\[
  J(\theta)=\E_{\tau\sim\pi_\theta}[R(\tau)].
\]

利用 log-derivative trick：

\[
  \nabla_\theta p_\theta(\tau)
  =
  p_\theta(\tau)\nabla_\theta\log p_\theta(\tau),
\]

得到 policy gradient 形式

\[
  \nabla_\theta J(\theta)
  =
  \E_{\pi_\theta}
  \left[
    \sum_{t\ge0}
    \nabla_\theta\log\pi_\theta(a_t\mid s_t)
    G_t
  \right].
\]

可用 \(Q^{\pi_\theta}(s_t,a_t)\) 或 advantage \(A^\pi(s_t,a_t)\) 替代 \(G_t\) 降低方差。baseline \(b(s)\) 不改变期望，因为

\[
  \E_{a\sim\pi_\theta(\cdot\mid s)}
  \left[
    \nabla_\theta\log\pi_\theta(a\mid s)b(s)
  \right]
  =
  b(s)\nabla_\theta\sum_a\pi_\theta(a\mid s)
  =
  0.
\]

Actor-critic 用 actor 表示策略，用 critic 估计值函数或 advantage。它把 policy gradient 的可优化性和 value learning 的低方差结合起来。

\paragraph{Policy gradient 的完整推导}

设一条轨迹为

\[
  \tau=(s_0,a_0,s_1,a_1,\ldots),
\]

轨迹回报为 \(R(\tau)\)。轨迹概率可分解为

\[
  p_\theta(\tau)
  =
  p(s_0)
  \prod_{t\ge0}
  \pi_\theta(a_t\mid s_t)
  P(s_{t+1}\mid s_t,a_t).
\]

环境转移 \(P(s_{t+1}\mid s_t,a_t)\) 不含 \(\theta\)，所以

\[
  \nabla_\theta\log p_\theta(\tau)
  =
  \sum_{t\ge0}
  \nabla_\theta\log\pi_\theta(a_t\mid s_t).
\]

目标函数为

\[
  J(\theta)
  =
  \int p_\theta(\tau)R(\tau)\,d\tau.
\]

对 \(\theta\) 求梯度：

\[
  \nabla_\theta J(\theta)
  =
  \int
  \nabla_\theta p_\theta(\tau)R(\tau)\,d\tau.
\]

使用 log-derivative trick，

\[
  \nabla_\theta p_\theta(\tau)
  =
  p_\theta(\tau)\nabla_\theta\log p_\theta(\tau),
\]

得到

\[
  \nabla_\theta J(\theta)
  =
  \E_{\tau\sim p_\theta}
  \left[
    R(\tau)
    \sum_{t\ge0}
    \nabla_\theta\log\pi_\theta(a_t\mid s_t)
  \right].
\]

进一步把整条轨迹回报替换为从 \(t\) 时刻开始的 return \(G_t\)，不会改变期望，但能减少无关过去奖励带来的方差，于是得到常用 policy gradient 形式。

\subsection{LLM 偏好对齐}

\concepttarget{concept:rlhf-dpo}

LLM 对齐可以看成特殊的序列决策问题。prompt 是初始状态，动作是逐 token 生成，策略是语言模型 \(\pi_\theta(o\mid x)\)，奖励来自人类或模型偏好。

RLHF 常先用偏好数据训练 reward model。若对同一 prompt \(x\) 有 preferred output \(o_w\) 和 rejected output \(o_l\)，Bradley-Terry 模型写作

\[
  P(o_w\succ o_l\mid x)
  =
  \sigma(r_\phi(x,o_w)-r_\phi(x,o_l)).
\]

奖励模型损失为

\[
  \Lcal_{\mathrm{RM}}
  =
  -
  \E
  \log
  \sigma(r_\phi(x,o_w)-r_\phi(x,o_l)).
\]

随后优化带 KL 正则的策略目标：

\[
  \max_{\pi_\theta}
  \E_{x,o\sim\pi_\theta}
  \left[
    r_\phi(x,o)
    -
    \beta
    \log
    \frac{\pi_\theta(o\mid x)}{\pi_{\mathrm{ref}}(o\mid x)}
  \right].
\]

KL 项防止模型为了 reward hacking 远离参考模型。DPO 则把奖励函数消去，直接用偏好对训练策略：

\[
  \Lcal_{\mathrm{DPO}}
  =
  -
  \E
  \log
  \sigma
  \left(
    \beta
    \log
    \frac{\pi_\theta(o_w\mid x)}{\pi_{\mathrm{ref}}(o_w\mid x)}
    -
    \beta
    \log
    \frac{\pi_\theta(o_l\mid x)}{\pi_{\mathrm{ref}}(o_l\mid x)}
  \right).
\]

考试若问“RLHF 和 DPO 差别”，重点是：RLHF 显式训练 reward model 并做策略优化；DPO 直接从偏好对构造监督式目标，不需要在线采样和独立 reward model。

\paragraph{DPO 目标的推导背景}

KL 正则 RLHF 的单个 prompt 目标可写成

\[
  \max_{\pi}
  \E_{o\sim\pi(\cdot\mid x)}
  \left[
    r(x,o)
    -
    \beta
    \log
    \frac{\pi(o\mid x)}{\pi_{\mathrm{ref}}(o\mid x)}
  \right].
\]

这个优化问题的闭式最优策略满足

\[
  \pi^*(o\mid x)
  =
  \frac{1}{Z(x)}
  \pi_{\mathrm{ref}}(o\mid x)
  \exp\left(\frac{1}{\beta}r(x,o)\right),
\]

其中 \(Z(x)\) 是归一化常数。移项得到

\[
  r(x,o)
  =
  \beta
  \log
  \frac{\pi^*(o\mid x)}{\pi_{\mathrm{ref}}(o\mid x)}
  +
  \beta\log Z(x).
\]

Bradley-Terry 偏好模型只使用奖励差：

\[
  r(x,o_w)-r(x,o_l).
\]

代入上式时，\(\beta\log Z(x)\) 在 winner 和 loser 之间相消，因此可用策略 log-ratio 直接表达偏好概率：

\[
  P(o_w\succ o_l\mid x)
  =
  \sigma
  \left(
    \beta
    \log
    \frac{\pi_\theta(o_w\mid x)}{\pi_{\mathrm{ref}}(o_w\mid x)}
    -
    \beta
    \log
    \frac{\pi_\theta(o_l\mid x)}{\pi_{\mathrm{ref}}(o_l\mid x)}
  \right).
\]

对这个概率取负对数似然，就是前面的 DPO loss \parencite{rafailov2023dpo,jurafsky2026slp}。

\paragraph{考核内容}

\begin{itemize}
  \item 定义 MDP、policy、return、\(V^\pi\)、\(Q^\pi\)。
  \item 从 return 递归推 Bellman expectation equation 和 optimality equation。
  \item 写出 Q-learning 更新式并解释 TD error。
  \item 用 log-derivative trick 推 policy gradient。
  \item 证明 baseline 不改变 policy gradient 期望。
  \item 解释 RLHF、KL 正则和 DPO 的数学目标及局限。
\end{itemize}

\section{自然语言处理与大模型理论}

\source{本节对应考纲“统计建模、Word2vec、self-attention、LLM 预训练对齐、Scaling Law、基于 Embedding 的知识库推理”。主要参考 \textcite{jurafsky2026slp,goldberg2017nnnlp}；Transformer、Scaling Law、RAG、DPO 等补充参考 \textcite{vaswani2017attention,kaplan2020scaling,hoffmann2022training,lewis2020rag,rafailov2023dpo}。}

\subsection{统计语言模型}

\concepttarget{concept:ngram}

语言模型给序列 \(w_{1:n}\) 分配概率。链式法则给出

\[
  P(w_{1:n})
  =
  \prod_{i=1}^{n}P(w_i\mid w_{1:i-1}).
\]

\(n\)-gram 模型用 Markov 假设截断历史：

\[
  P(w_i\mid w_{1:i-1})
  \approx
  P(w_i\mid w_{i-n+1:i-1}).
\]

最大似然估计为

\[
  \widehat P(w_i\mid c)
  =
  \frac{\operatorname{count}(c,w_i)}
  {\operatorname{count}(c)}.
\]

问题是稀疏性：很多合理短语在训练集中没出现，MLE 会给零概率。因此需要 smoothing、backoff 或 interpolation。困惑度（perplexity）衡量测试序列平均预测难度：

\[
  \operatorname{PPL}(w_{1:n})
  =
  P(w_{1:n})^{-1/n}.
\]

PPL 越低，模型对测试文本赋予的平均概率越高。

\paragraph{\(n\)-gram MLE 的推导}

固定一个上下文 \(c\)。设词表为 \(\mathcal{V}\)，参数为

\[
  \theta_w=P(w\mid c),
  \qquad
  \sum_{w\in\mathcal{V}}\theta_w=1.
\]

给定计数 \(n(c,w)=\operatorname{count}(c,w)\)，该上下文相关的 log-likelihood 为

\[
  \sum_{w\in\mathcal{V}}
  n(c,w)\log\theta_w.
\]

引入拉格朗日乘子 \(\lambda\)：

\[
  \mathcal{J}(\theta,\lambda)
  =
  \sum_w n(c,w)\log\theta_w
  +
  \lambda\left(\sum_w\theta_w-1\right).
\]

一阶条件为

\[
  \frac{n(c,w)}{\theta_w}+\lambda=0,
  \qquad
  \theta_w=-\frac{n(c,w)}{\lambda}.
\]

对 \(w\) 求和并使用归一化约束：

\[
  1=\sum_w\theta_w
  =
  -\frac{1}{\lambda}
  \sum_w n(c,w)
  =
  -\frac{\operatorname{count}(c)}{\lambda}.
\]

所以 \(\lambda=-\operatorname{count}(c)\)，得到

\[
  \widehat P(w\mid c)
  =
  \frac{\operatorname{count}(c,w)}{\operatorname{count}(c)}.
\]

\subsection{Word2vec、GloVe 与 embedding}

\concepttarget{concept:word2vec}

Embedding 把离散词、实体或节点映射到向量空间：

\[
  e:\mathcal{V}\to\R^d.
\]

Word2vec skip-gram 用中心词预测上下文。给中心词 \(w\) 和上下文词 \(c\)，softmax 概率为

\[
  P(c\mid w)
  =
  \frac{\exp(u_c^\top v_w)}
  {\sum_{c'}\exp(u_{c'}^\top v_w)}.
\]

完整 softmax 太贵，negative sampling 用二分类目标近似：

\[
  \log\sigma(u_c^\top v_w)
  +
  \sum_{j=1}^{k}
  \E_{c_j\sim P_n}
  \log\sigma(-u_{c_j}^\top v_w).
\]

GloVe 从全局共现矩阵出发，使

\[
  u_i^\top v_j+b_i+\tilde b_j
  \approx
  \log X_{ij}.
\]

二者共同直觉是 distributional hypothesis：上下文相似的词语语义相似。

在 skip-gram 公式中，\(v_w\in\R^d\) 是中心词 \(w\) 的输入向量，\(u_c\in\R^d\) 是上下文词 \(c\) 的输出向量，\(P_n\) 是负采样分布，\(k\) 是每个正样本配多少个负样本。考试如果要求解释 negative sampling，应说明它把多分类 softmax 近似成“真实共现对 vs. 噪声共现对”的二分类问题，从而避免每次更新遍历整个词表 \(\mathcal{V}\)。

\subsection{Transformer self-attention}

\concepttarget{concept:self-attention}

Transformer 的输入是 token embedding 矩阵

\[
  X\in\R^{N\times d},
\]

其中 \(N\) 是上下文长度，\(d\) 是模型维度。单头 attention 先计算

\[
  Q=XW_Q,
  \qquad
  K=XW_K,
  \qquad
  V_{\mathrm{attn}}=XW_V.
\]

然后

\[
  \operatorname{Attention}(Q,K,V_{\mathrm{attn}})
  =
  \softmax
  \left(
    \frac{QK^\top}{\sqrt{d_k}}
  \right)V_{\mathrm{attn}}.
\]

每个 query 与所有 key 做相似度，softmax 得到权重，再对 value 加权平均。缩放因子 \(1/\sqrt{d_k}\) 防止点积方差随维度增大，使 softmax 过早饱和。

\paragraph{Attention 维度和缩放因子的推导}

设 \(W_Q,W_K\in\R^{d\times d_k}\)，\(W_V\in\R^{d\times d_v}\)。则

\[
  Q,K\in\R^{N\times d_k},
  \qquad
  V_{\mathrm{attn}}\in\R^{N\times d_v}.
\]

所以

\[
  QK^\top\in\R^{N\times N},
  \qquad
  \softmax(QK^\top/\sqrt{d_k})V_{\mathrm{attn}}\in\R^{N\times d_v}.
\]

若 query 和 key 的各坐标近似独立、均值为 \(0\)、方差为 \(1\)，则一个点积

\[
  q^\top k=\sum_{r=1}^{d_k}q_r k_r
\]

的方差为

\[
  \Var(q^\top k)
  =
  \sum_{r=1}^{d_k}\Var(q_r k_r)
  =
  d_k.
\]

因此未缩放的 logits 会随 \(d_k\) 变大，softmax 进入饱和区后梯度变小。除以 \(\sqrt{d_k}\) 后，点积方差回到常数量级，这是 scaled dot-product attention 的动机 \parencite{vaswani2017attention,jurafsky2026slp}。

自回归语言模型必须加 causal mask：

\[
  M_{ij}
  =
  \begin{cases}
    0, & j\le i,\\
    -\infty, & j>i.
  \end{cases}
\]

此时 attention 变成

\[
  \softmax
  \left(
    \frac{QK^\top}{\sqrt{d_k}}+M
  \right)V_{\mathrm{attn}},
\]

保证第 \(i\) 个位置不能偷看未来 token。这里 \(M\) 是加到 logits 上的 causal mask；padding mask 则遮蔽序列填充位置，二者不可混用。实现时还要核对标签是否已右移：若预测第 \(i\) 个目标 token，严格下三角掩码与输入/标签的移位约定必须匹配。多头 attention 把多个 head 的输出拼接后再线性投影。Transformer block 通常包括 attention、feed-forward network、residual connection 和 layer norm。

\subsection{LLM 预训练、指令微调与对齐}

\concepttarget{concept:pretraining-alignment}

自回归 LLM 的预训练目标是 next-token prediction：

\[
  \Lcal_{\mathrm{pretrain}}
  =
  -
  \sum_{i=1}^{n}
  \log p_\theta(w_i\mid w_{<i}).
\]

这就是交叉熵损失。它使模型学习语言统计、事实模式和许多任务的隐式格式，但目标本身只是预测文本，不保证有帮助、真实或安全。

这个目标同样来自最大似然。给定语料序列 \(w_{1:n}\)，自回归分解为

\[
  p_\theta(w_{1:n})
  =
  \prod_{i=1}^{n}p_\theta(w_i\mid w_{<i}).
\]

最大化 log-likelihood 等价于最小化

\[
  -\log p_\theta(w_{1:n})
  =
  -
  \sum_{i=1}^{n}
  \log p_\theta(w_i\mid w_{<i}).
\]

如果把真实下一个 token 写成 one-hot 分布 \(q_i\)，模型预测分布写成 \(p_i\)，单 token 损失就是交叉熵

\[
  H(q_i,p_i)
  =
  -
  \sum_{v\in\mathcal{V}}q_i(v)\log p_i(v)
  =
  -\log p_\theta(w_i\mid w_{<i}).
\]

指令微调（SFT）继续用交叉熵训练，但数据变成 instruction-response pairs：

\[
  \Lcal_{\mathrm{SFT}}
  =
  -
  \sum_{(x,y)}
  \sum_{t}
  \log p_\theta(y_t\mid x,y_{<t}).
\]

偏好对齐在上一节已给出 RLHF 与 DPO。它的局限包括 reward model 偏差、偏好数据覆盖不足、过度优化导致 reward hacking、长程事实性仍依赖检索或工具。

\subsection{Scaling Law}

\concepttarget{concept:scaling-law}

Scaling law 研究模型性能如何随参数量 \(N\)、数据 token 数 \(D\)、训练计算量 \(C\) 缩放。经验上，交叉熵损失常近似满足幂律：

\[
  L(N,D)
  \approx
  L_\infty
  +
  aN^{-\alpha}
  +
  bD^{-\beta}.
\]

这里 \(L_\infty\) 是不可约损失水平，\(a,b>0\) 是经验拟合常数，\(\alpha,\beta>0\) 是参数量和数据量对应的幂律指数。

这不是一个第一性原理定理，而是大规模实验规律。它的考试价值在于解释资源分配：固定 compute 时，参数量和数据量要平衡；只增大模型但数据不足会 data-limited，只增加数据但模型太小会 model-limited。Chinchilla 系列结果强调 compute-optimal 训练往往需要比早期大模型更多 token、更少参数。

\subsection{Embedding 知识库推理、RAG 与 MLN}

\concepttarget{concept:rag-mln}

知识库 embedding 把实体和关系映射到向量空间。TransE 的基本约束是

\[
  e_s+r\approx e_o
\]

对三元组 \((s,r,o)\) 成立。评分函数可写为

\[
  f(s,r,o)=-\|e_s+r-e_o\|.
\]

局限是多对多关系、组合逻辑、否定和量词难以表达。图神经网络、规则增强、路径推理和神经符号方法都是常见改进方向。

RAG 把检索器和生成器组合。给定查询 \(q\)，检索器从文档库 \(D\) 返回

\[
  R(q)=\{d_1,\ldots,d_k\},
\]

生成器根据 \(q\) 和检索文档生成答案：

\[
  p(x_{1:n}\mid q,R(q))
  =
  \prod_{i=1}^{n}
  p(x_i\mid q,R(q),x_{<i}).
\]

RAG 的价值是缓解幻觉、接入私有或时效知识、提供引用证据。主要风险是检索失败、检索噪声、上下文污染、重排错误和引用不忠实。工程闭环通常包括 chunking、embedding、向量索引、召回、重排、生成、引用和评估。

Markov Logic Network（MLN）用加权一阶逻辑公式定义 log-linear 分布：

\[
  P(X=x\mid E)
  =
  \frac{1}{Z(E)}
  \exp
  \left(
    \sum_j w_j n_j(x,E)
  \right).
\]

其中 \(n_j(x,E)\) 是第 \(j\) 个公式在世界 \(x\) 与证据 \(E\) 下满足的 groundings 数，\(w_j\) 是该公式的权重，\(Z(E)=\sum_x \exp(\sum_j w_j n_j(x,E))\) 是在证据 \(E\) 下的配分函数，用来保证概率和为 \(1\)。权重大表示强偏好；无限大权重可视为 hard constraint。它和 embedding 方法的差别是：MLN 显式编码逻辑约束，embedding 方法更擅长连续相似性和大规模近似。

\paragraph{考核内容}

\begin{itemize}
  \item 从链式法则写出语言模型目标，解释 \(n\)-gram 的 Markov 假设与 smoothing。
  \item 推 Word2vec skip-gram 或 negative sampling 目标，解释 embedding 的语义来源。
  \item 写出 self-attention 的 \(Q,K,V\) 公式、维度、mask 和复杂度。
  \item 解释 next-token pretraining、SFT、RLHF、DPO 的目标差异。
  \item 说明 scaling law 的经验形式和 compute-optimal trade-off。
  \item 比较知识库 embedding、RAG 和 MLN 的能力边界。
\end{itemize}

\section{第三章自测：定义、条件与反例}

以下问题用于检查是否真正掌握了“对象--假设--结论”链条。每题应先写定义，再列出所需条件，最后给出一行反例或失败模式；这样写出的答案也更接近资格考试的评分点。

\begin{enumerate}
  \item \textbf{泛化与模型选择。} 为什么有限类的 union bound 不能直接替代 VC 证明？请写出 ghost sample 的作用；再区分类内最优风险、近似误差和近似 ERM 的优化误差。
  \item \textbf{经典模型。} 在线性回归中何时正规方程不可逆？岭回归怎样改变该问题？给出一个在高维未缩放特征下 \(k\)-NN 距离失真的例子。
  \item \textbf{核与压缩感知。} 为什么“某个 Gram 矩阵半正定”不足以定义核？\(2s\)-RIP 保证的唯一性是什么，为什么它仍不足以单独证明 \(\ell_1\) 恢复？
  \item \textbf{优化。} 强凸性与极小点存在性分别提供什么？请比较 GD、SGD 和随机坐标下降的随机来源，并列出 Adam 的一阶矩、二阶矩、偏差修正与数值稳定项。
  \item \textbf{深度学习。} 万能逼近定理没有保证哪两件事？任选视觉任务，写出其输入、输出、训练损失与一个常见失效模式。
  \item \textbf{强化学习。} Bellman 最优算子在哪个范数下是压缩映射？表格型 Q-learning 的收敛需要哪些访问、步长与噪声条件？
  \item \textbf{NLP/LLM。} 写出 attention logits、causal mask 与 padding mask 的区别；再说明预训练、SFT、偏好优化分别改变什么目标，以及 RAG 不能自动消除的失败模式。
\end{enumerate}

\section{从第三章到真题的使用方法}

第三章不是单独背诵材料，而是后续真题解答的工具箱。遇到题目时，先判断它属于哪门核心课，再调用对应的定义和推导模板。

\begin{longtable}{L{0.24\textwidth}L{0.34\textwidth}L{0.34\textwidth}}
\toprule
考纲主题 & 必备工具 & 常见考法 \\
\midrule
\endfirsthead
\toprule
考纲主题 & 必备工具 & 常见考法 \\
\midrule
\endhead
PAC/VC & Bernoulli 指示变量、Hoeffding、union bound、Sauer 引理 & 推样本复杂度，证明统一收敛，计算或界定 VC 维 \\
模型选择 & approximation/estimation error、bias-variance、validation & 解释欠拟合/过拟合，设计正则化或交叉验证方案 \\
线性/逻辑模型 & 正规方程、MLE、logistic loss、梯度 & 推目标函数、梯度、凸性和正则化效果 \\
SVM/核 & margin、hinge loss、KKT 直觉、PDS kernel & 写 hard/soft margin，解释核技巧和 margin 泛化 \\
压缩感知 & 稀疏性、欠定测量、\(\ell_1\) 松弛、RIP & 解释为什么少量随机测量可恢复稀疏信号 \\
凸优化 & 一阶条件、smoothness、强凸性、次梯度、prox & 证明下降、收敛率，推 prox 或 SGD 递推 \\
深度学习 & 反向传播、归一化、残差、生成模型目标 & 推梯度，比较 VAE/GAN/flow/diffusion，解释训练稳定性 \\
强化学习 & MDP、Bellman、TD error、policy gradient & 推 Bellman 方程、Q-learning、baseline 无偏性 \\
NLP/LLM & n-gram、embedding、attention、pretraining、alignment、RAG & 写 attention 维度，解释 SFT/RLHF/DPO、RAG 失败模式 \\
\bottomrule
\end{longtable}

\begin{checkpoint}
第三章的复习顺序建议是：先掌握风险与泛化，再掌握凸优化；然后学经典模型和深度网络；最后学 RL 与 NLP。原因是后两者的大多数公式都在复用前两部分的概率、优化和函数逼近语言。
\end{checkpoint}

\examnote
做真题时可以按以下模板组织答案：

\begin{enumerate}
  \item 写清对象：空间、随机变量、参数、损失或值函数。
  \item 写目标：经验风险、期望回报、似然、ELBO、Bellman residual 或对齐目标。
  \item 写关键性质：凸性、无偏性、集中界、递归关系、可逆性、mask 结构。
  \item 写结论：样本复杂度、梯度、更新式、泛化解释或模型局限。
  \item 若是开放题，补工程风险：数据泄漏、分布漂移、检索失败、优化不稳定、计算成本。
\end{enumerate}
